%%%%%%%% ICML 2026 EXAMPLE LATEX SUBMISSION FILE %%%%%%%%%%%%%%%%%

\documentclass{article}

% Recommended, but optional, packages for figures and better typesetting:
\usepackage{microtype}
\usepackage{graphicx}
\usepackage{subcaption}
\usepackage{booktabs} % for professional tables

% hyperref makes hyperlinks in the resulting PDF.
% If your build breaks (sometimes temporarily if a hyperlink spans a page)
% please comment out the following usepackage line and replace
% \usepackage{icml2026} with \usepackage[nohyperref]{icml2026} above.
\usepackage{hyperref}


% Attempt to make hyperref and algorithmic work together better:
\newcommand{\theHalgorithm}{\arabic{algorithm}}

% Use the following line for the initial blind version submitted for review:
% \usepackage{icml2026}

% For preprint, use
\usepackage[preprint]{icml2026}

% If accepted, instead use the following line for the camera-ready submission:
% \usepackage[accepted]{icml2026}

\usepackage{amsmath}
\usepackage{amssymb}
\usepackage{mathtools}
\usepackage{amsthm}


% #### custom
\usepackage{algorithm}
\usepackage{algorithmic}

\usepackage{tabularx}


% if you use cleveref..
\usepackage[capitalize,noabbrev]{cleveref}

%%%%%%%%%%%%%%%%%%%%%%%%%%%%%%%%
% THEOREMS
%%%%%%%%%%%%%%%%%%%%%%%%%%%%%%%%
\theoremstyle{plain}
\newtheorem{theorem}{Theorem}[section]
\newtheorem{proposition}[theorem]{Proposition}
\newtheorem{lemma}[theorem]{Lemma}
\newtheorem{corollary}[theorem]{Corollary}
\theoremstyle{definition}
\newtheorem{definition}[theorem]{Definition}
\newtheorem{assumption}[theorem]{Assumption}
\theoremstyle{remark}
\newtheorem{remark}[theorem]{Remark}

% Todonotes is useful during development; simply uncomment the next line
%    and comment out the line below the next line to turn off comments
%\usepackage[disable,textsize=tiny]{todonotes}

% The \icmltitle you define below is probably too long as a header.
% Therefore, a short form for the running title is supplied here:
\icmltitlerunning{Self-Conditioned Denoising for Atomistic Representation Learning}

\begin{document}

\twocolumn[
  \icmltitle{Self-Conditioned Denoising for Atomistic Representation Learning}

  % It is OKAY to include author information, even for blind submissions: the
  % style file will automatically remove it for you unless you've provided
  % the [accepted] option to the icml2026 package.

  % List of affiliations: The first argument should be a (short) identifier you
  % will use later to specify author affiliations Academic affiliations
  % should list Department, University, City, Region, Country Industry
  % affiliations should list Company, City, Region, Country

  % You can specify symbols, otherwise they are numbered in order. Ideally, you
  % should not use this facility. Affiliations will be numbered in order of
  % appearance and this is the preferred way.
  \icmlsetsymbol{equal}{*}

  \begin{icmlauthorlist}
    \icmlauthor{Tynan J. Perez}{yyy}
    \icmlauthor{Rafael Gómez-Bombarelli}{xxx}
    % \icmlauthor{Firstname3 Lastname3}{comp}
    % \icmlauthor{Firstname4 Lastname4}{sch}
    % \icmlauthor{Firstname5 Lastname5}{yyy}
    % \icmlauthor{Firstname6 Lastname6}{sch,yyy,comp}
    % \icmlauthor{Firstname7 Lastname7}{comp}
    % %\icmlauthor{}{sch}
    % \icmlauthor{Firstname8 Lastname8}{sch}
    % \icmlauthor{Firstname8 Lastname8}{yyy,comp}
    % %\icmlauthor{}{sch}
    % %\icmlauthor{}{sch}
  \end{icmlauthorlist}

  \icmlaffiliation{yyy}{Department of Chemistry, Massachusetts
Institute of Technology, 77 Massachusetts Avenue, Cambridge, 02139,
MA, USA}
  \icmlaffiliation{xxx}{Department of Materials Science and Engineering, Massachusetts
Institute of Technology, 77 Massachusetts Avenue, Cambridge, 02139,
MA, USA}
  % \icmlaffiliation{comp}{Company Name, Location, Country}
  % \icmlaffiliation{sch}{School of ZZZ, Institute of WWW, Location, Country}

  \icmlcorrespondingauthor{Rafael Gómez-Bombarelli}{rafagb@mit.edu}

  % You may provide any keywords that you find helpful for describing your
  % paper; these are used to populate the "keywords" metadata in the PDF but
  % will not be shown in the document
  \icmlkeywords{Machine Learning, ICML}

  \vskip 0.3in
]

% this must go after the closing bracket ] following \twocolumn[ ...

% This command actually creates the footnote in the first column listing the
% affiliations and the copyright notice. The command takes one argument, which
% is text to display at the start of the footnote. The \icmlEqualContribution
% command is standard text for equal contribution. Remove it (just {}) if you
% do not need this facility.

% Use ONE of the following lines. DO NOT remove the command.
% If you have no special notice, KEEP empty braces:

\printAffiliationsAndNotice{}  % no special notice (required even if empty)

% Or, if applicable, use the standard equal contribution text:
% \printAffiliationsAndNotice{\icmlEqualContribution}

\begin{abstract}
  % This document provides a basic paper template and submission guidelines.
  % Abstracts must be a single paragraph, ideally between 4--6 sentences long.
  % Gross violations will trigger corrections at the camera-ready phase.
  

The success of large-scale pretraining in NLP and computer vision has catalyzed growing efforts to develop analogous foundation models for the physical sciences. However,  pretraining strategies using atomistic data remain underexplored. To date, large-scale supervised pretraining on DFT force-energy labels has provided the strongest performance gains to downstream property prediction, out-performing existing methods of self-supervised learning (SSL) which remain limited to ground-state geometries, and/or single domains of atomistic data. We address these shortcomings with Self-Conditioned Denoising (SCD), a backbone-agnostic reconstruction objective that utilizes self-embeddings for conditional denoising across any domain of atomistic data, including small molecules, proteins, periodic materials, and 'non-equilibrium' geometries. When controlled for backbone architecture and pretraining dataset, SCD significantly outperforms previous SSL methods on downstream benchmarks and matches or exceeds the performance of supervised force-energy pretraining. We show that a small, fast GNN pretrained by SCD can achieve competitive or superior performance to larger models pretrained on significantly larger labeled or unlabeled datasets, across tasks in multiple domains. Our code is available at: \url{https://github.com/TyJPerez/SelfConditionedDenoisingAtoms}

\end{abstract}

\section{Introduction}

The development of 'foundation models' has transformed natural language processing (NLP) and computer vision. The prevailing training paradigm now employs a two-stage process: pretraining on large unlabeled datasets to learn generalizable representations, followed by fine-tuning on smaller, labeled datasets for specialized tasks. The success of large-scale pretraining has motivated the development of analogous foundation models in atomistic sciences, including chemistry, structural biology, and materials science \cite{AtomisticfFoundationModels_yuan2025, AI4sci_review_zhang2025artificial}. These efforts are driven by their potential for accelerated atomistic simulation, direct structure-to-property prediction, and inverse design—capabilities that could significantly accelerate the discovery of new drugs and materials.

In chemistry and materials science, the most notable progress toward this vision has come from researchers developing machine-learned interatomic potentials (MLIPs) to accelerate atomistic simulation. Pretrained MLIPs - such as UMA\cite{UMA_wood2025}, MACE\cite{MACE_batatia2025} and Orb\cite{Orbv3_rhodes2025} - have achieved remarkable accuracy in DFT force-energy predictions by supervised pretraining on large, diverse atomistic datasets. These models have enabled near-DFT quality simulations and structural relaxations at speeds and scales far exceeding those of traditional DFT. While often termed 'foundation models' (or more appropriately 'foundation potentials'), they differ from their counterparts in NLP and CV in one key regard: current foundation potentials are pretrained exclusively via supervised learning on precomputed forces and energies. However, scaling supervised pretraining with DFT labels is computationally expensive; generating forces and energies for the 100M geometries in the OMol25 dataset required six billion CPU core-hours \cite{OMOL25_levine2025open}.

In other fields, foundation models typically rely on self-supervised learning (SSL) objectives that leverage vast, unlabeled datasets. However, unlike images and text, which are often collected without labels, atomistic datasets are typically generated via computational methods (such as DFT) that inherently produce force and energy labels as part of the data generation process. Although force and energy values vary depending on the DFT functional used, prior work has overcome this by using multiple prediction heads and demonstrating clear benefits to supervised pretraining on large, diverse datasets containing labels derived from multiple functionals \cite{UMA_wood2025, JMP_Shoghi2023}. Consequently, supervised pretraining with DFT labels remains the dominant paradigm, outperforming existing SSL methods on simulation and property prediction benchmarks \cite{JMP_Shoghi2023, NoisyNodes_Zaidi2023PVD}. However, the ratio of 'ground-truth' DFT labeled atomistic data is likely to change as more atomistic geometries are produced using MLIP-driven simulation or relaxation, and generative models \cite{boltz1_2024, boltz2_2025, Sair_dataset, UMA_wood2025, mattergen_zeni2023, MACE_batatia2025, GeomRep_li2024}. Nevertheless, adapting SSL methods from other fields has remained a challenge due to the unique attributes of 3D atomistic data. Atomistic data has no fixed size (samples range from 3 to over 3000 atoms), samples may exhibit periodicity in 3D space, and learned representations must be sensitive to small changes in atomic position to achieve accurate predictions in geometrically sensitive tasks (structural relaxation, simulation, binding affinity predictions, etc). 

Previous work in computer vision suggests that self-supervised pretraining can yield more interpretable and generalizable representations than supervised learning \cite{Mocov3_chen2021, Dino2_oquab2024, DINO3_siméoni2025dinov3}, achieving performance competitive with supervised models even on identical datasets \cite{MAE_he2021, ReLicv2_tomasev2022}. This is because supervised learning only requires a model to represent features of the training data that are relevant to a specific task, becoming invariant to uncorrelated features. In contrast, generative SSL encourages a model to represent the full feature distribution within a training dataset, often producing internal representations with superior transferability. In this work, we demonstrate that this principle holds for 3D atomistic data as well.


% Crucially, atomistic data pose unique challenges when adapting SSL methods from other fields.

% In other fields, foundation models typically rely on self-supervised learning (SSL) objectives that leverage vast, unlabeled datasets. However, unlike images and text, 
% atomistic data poses challenges for SSL. Atomistic data has no fixed size (samples range from 3 to over 3000 atoms), samples may exhibit periodicity in 3D space, and learned representations must be sensitive to small changes in atomic position to achieve accurate predictions in geometrically sensitive tasks (structural relaxation, simulation, binding affinity predictions, etc). 
% Further, compared to language or vision data, a large fraction of currently available atomistic geometries have corresponding DFT-calculated force and energy labels. 
% Unlike images or textual data, which are often collected without labels, atomistic datasets are typically generated via computational methods (such as DFT) that inherently produce force and energy labels as part of the data generation process.
% Although force and energy values vary depending on the DFT functional used, prior work has overcome this by using multiple prediction heads, demonstrating clear benefits to supervised pretraining on large datasets containing labels derived from multiple functionals \cite{UMA_wood2025, JMP_Shoghi2023}. As a result, supervised pretraining with DFT labels remains the dominant approach, currently yielding superior performance for both simulation and property prediction. However, the ratio of DFT labeled atomistic data is likely to change as more atomistic geometries are produced using MLIP-driven simulation or relaxation, and generative models \cite{boltz1_2024, boltz2_2025, Sair_dataset, UMA_wood2025, mattergen_zeni2023, MACE_batatia2025, GeomRep_li2024}.

\section{Background}
\subsection{Contrastive Learning on Low Dimensional Representation}
Early atomistic SSL relied on maximizing mutual information between 3D geometries and lower-dimensional representations, such as 2D graphs or SMILES \cite{GraphMVP_liu2022, 3dinfomax_stärk2022, attrmasking_hu2020}. While these methods improved performance for predictions on lower dimensional inputs, they often underperformed supervised models acting on 3D inputs, and were not applicable in configurationally aware tasks like force prediction or ligand docking. Furthermore, multi-modal contrastive approaches are inapplicable to systems lacking 1D representations, such as periodic crystals, limiting their generality.

\subsection{Node Denoising}
Node denoising pretraining emerged as a more generalizable alternative, operating directly on 3D point clouds. Formally, given an input structure with N atoms $\mathbf{x} \in \mathbb{R}^{N\times3}$, we sample Gaussian noise $\boldsymbol{\varepsilon} \sim \mathcal{N}(\mathbf{0}, \sigma^{2}\mathbf{I})$ to generate a corrupted geometry $\tilde{\mathbf{x}} = \mathbf{x} + \boldsymbol{\varepsilon}$. A neural network $\phi_{\theta}$ is then trained to predict the noise vector on each atom of the corrupted geometry:
\begin{equation}
\label{eq:denoising}
\mathcal{L}_{\text{denoise}} = \mathbb{E}_{\mathbf{x}, \boldsymbol{\varepsilon}} \left[ \left\lVert \phi_{\theta}(\tilde{\mathbf{x}}) - \boldsymbol{\varepsilon} \right\rVert_{2}^{2} \right]
\end{equation}
Unlike contrastive objectives that promote multi-view invariance, noise prediction encourages equivariance to geometric perturbations, making it effective for learning geometric features of local atomic environments. Pretraining by node denoising  has proven useful as both an auxiliary loss \cite{NoisyNodes_Godwin2022, Equiformer1_Liao2022, EquiformerV2_Liao2023} and a standalone pretraining task \cite{NoisyNodes_Zaidi2023PVD, SE3invDenoise_liu2023, Atomica_Fang2025, EPT_jiao2025}.

% \subsection{Variations on Node Denoising}
\subsection{State of The Art Pretraining}
Despite its elegance and generality, standard node denoising struggles to match the efficacy of supervised force-energy (FE) pretraining, largely because it fails to benefit from training on non-ground state ('non-equilibrium') geometries \cite{ETOREO_feng2023}. Early interpretations suggested that denoising approximates learning a 'near-equilibrium' force field \cite{NoisyNodes_Zaidi2023PVD}. However, for high-energy configurations, noise vectors $\boldsymbol{\varepsilon}$ (the vectors pointing from corrupted atom positions to uncorrupted positions) are unlikely to resemble true forces and the difference between a corrupted input and a valid non-corrupted higher-energy configuration becomes ambiguous. This makes denoising 'non-equilibrium' geometries an ill-posed task without additional information. One solution is to incorporate a force encoding \cite{DeNS_liao2024}, but this is no longer SSL as it requires force-energy labels.  

Others have shown that there is little correlation between noise vectors and true DFT forces even for ground-state geometries \cite{SliDe_Ni2025}.  Variations on the node denoising objective have improved performance, but at the cost of introducing domain-specific constraints, such as torsional noise in \textit{Frad} \cite{Frad_Ni2024}, or pseudo-force targets in \textit{SliDe} \cite{SliDe_Ni2025}. Others have explored specialized masking, block-wise denoising,  hybrid contrastive objectives with lower-dimensional representations, and scaling pretraining on larger datasets with to up to 200M geometries \cite{Atomica_Fang2025,UniCorn_Feng2024,EPT_jiao2025, UniMol_Zhou2023}. However, these approaches often sacrifice generality and remain limited to ground state geometries of small molecules. To date, general-purpose SSL objectives have not matched the performance of supervised FE pretraining on downstream property prediction tasks \cite{JMP_Shoghi2023, HackNIP_kim2025}.



\section{Method}
\subsection{The Limitations of Node Denoising}

We identify three critical limitations of standard node denoising as a pretraining objective:

\begin{description}
\item[\textbf{(s1) Locality of Node Denoising:}] 
GNNs often struggle with global representations and long-range correlations. Node denoising fails to produce strong global semantic learning because: (1) Gaussian noise applied to atom positions lacks long-range correlation across space, and (2) small perturbations in atom position can be estimated with high accuracy using a relatively small, local context window.

\item[\textbf{(s2) Scalar vs. Vector Embedding Pressure:}] 
 Equivariant GNNs commonly partition node embeddings into multiple channels determined by rotational symmetry groups ($L=0, 1, 2, \dots$). Because noise prediction is an equivariant vector ($L=1$) objective, it primarily impacts $L=1$ channels. While denoising gradients propagate through all model parameters, there is insufficient pressure on invariant scalar embeddings ($L=0$) to develop good semantic representations. This is problematic for downstream property predictions, which are often roto-translation invariant and rely on the invariant ($L=0$) channel.

\item[\textbf{(s3) Ambiguity with Non-Equilibrium Structures:}] 
Standard node denoising pretraining does not benefit from including 'non-equilibrium' or high-energy geometries. Without additional context,  models cannot distinguish between a corrupted input geometry and a valid, high-energy target geometry. This creates ambiguity in the objective when training on datasets that include non-ground-state conformers.
\end{description}

\subsection{Self-Conditioned Denoising}
Self-Conditioned Denoising (SCD) addresses all three shortcomings of the standard denoising objective by incorporating a conditional embedding in the model architecture as shown in eqn. \ref{scd_obj1}:
\begin{equation}\label{scd_obj1}
\mathbb{E}_{q_{\sigma}(\tilde{\mathbf{x}}, \mathbf{x})}
\left[
  \left\lVert
    \phi_{\theta}(\tilde{\mathbf{x}}|\mathbf{c}) - \varepsilon
  \right\rVert_{2}^{2}
\right], \quad
\mathbf{c} = f_{\eta}(\mathbf{x})
\end{equation}
Here, $c \in  \mathbb{R}^{d}$ is an embedding vector that represents the uncorrupted target geometry. This conditioning provides the model with enough information to distinguish between corrupted samples and a valid target conformation (addressing s3). The conditional embedding itself can then be used as a training objective to develop meaningful invariant (L0) representation (addressing s2). This in turn pushes the embedding model ($f_{\eta}$) to develop better global/semantic representation (addressing s1). In SCD, we use the same model as both the embedding model $f_{\eta}$, and the noise predictor, $\phi_{\theta}$, during each optimizer step of pretraining (shown in eqn \ref{SCD_obj2}).
\begin{equation}\label{SCD_obj2}
\begin{aligned}
\mathbb{E}_{q_{\sigma}(\tilde{\mathbf{x}}, \mathbf{x})}
\left[
  \left\lVert
    \phi_{\theta}(\tilde{\mathbf{x}}|\mathbf{c}) - \varepsilon
  \right\rVert_{2}^{2}
\right]\\[1ex]
=
\mathbb{E}_{q_{\sigma}(\tilde{\mathbf{x}}, \mathbf{x})}
\left[
  \left\lVert
    \phi_{\theta}(\tilde{\mathbf{x}}|\phi_{\theta}(\mathbf{x})_{L0})_{L1} - \varepsilon
  \right\rVert_{2}^{2}
\right]\\[2ex]
\mathbf{c} = \phi_{\theta}(\mathbf{x})_{L0} \in \mathbb{R}^d, \quad
\phi_{\theta}(\tilde{\mathbf{x}}|\mathbf{c})_{L1} \in \mathbb{R}^{N\times3}\\
\end{aligned}
\end{equation}
Implementation of the SCD objective is summarized by algorithm \ref{alg:pretrain_finetune} and visualized in figure \ref{fig:scd}. First a forward pass is made using uncorrupted geometries ($\mathbf{x}$). The resulting invariant (L0) embeddings are sum-pooled across nodes and passed through a two-layer MLP embedding head, forming a single embedding per geometry. This creates an information bottleneck that encourages semantic representation. The embedding is then used as conditional guidance for a second forward pass during which a corrupted geometry ($\tilde{\mathbf{x}}$) is used as input. The L1 equivariant output from the second pass is then used to predict the per-node noise corruption. To preserve unconditional behavior, we drop 20\% of conditional embeddings each pass. After pretraining is complete, downstream supervised finetuning and inference only requires a single model pass without conditioning. 


% ########## DOUBLE COLUMN VERSION ######
% \begin{algorithm*}[t]
% \small
% \caption{Pretraining and Fine-tuning Pipeline}
% \label{alg:pretrain_finetune}
% \begin{algorithmic}[1]
% \REQUIRE Dataset $\mathcal{D}$, noise scales $\sigma_{\text{corr}}, \sigma_{\text{reg}}$, model $\phi_{\theta}$
% \ENSURE Trained parameters $\theta$

% \vspace{0.5em}
% \STATE \textbf{Pretraining}
% \FOR{each minibatch $\mathbf{x} \sim \mathcal{D}$}
%     \STATE $(\mathbf{x}, \tilde{\mathbf{x}}, \varepsilon) \leftarrow \textsc{Corrupt}(\mathbf{x}, \sigma_{\text{corr}}, \sigma_{\text{reg}})$ \hfill $\triangleright$ clean and corrupted views
%     \STATE $(\mathbf{c}_{L0}, \mathbf{v}_{L1}, \mathbf{y}_{L0}) \leftarrow \phi_{\theta}(\mathbf{x})$
%     \STATE $(\tilde{\mathbf{c}}_{L0}, \tilde{\mathbf{v}}_{L1}, \tilde{\mathbf{y}}_{L0}) \leftarrow \phi_{\theta}(\tilde{\mathbf{x}}, \mathbf{c}_{L0})$
%     \STATE $\mathcal{L}_{\text{pre}} \leftarrow \mathrm{MSE}(\tilde{\mathbf{v}}_{L1}, \varepsilon)$
%     \STATE $\theta \leftarrow \theta - \eta \nabla_{\theta} \mathcal{L}_{\text{pre}}$
% \ENDFOR

% \vspace{0.5em}
% \STATE \textbf{Fine-tuning / Inference}
% \FOR{each labeled sample $(\mathbf{x}, \mathbf{y})$}
%     \STATE $(\mathbf{c}_{L0}, \mathbf{v}_{L1}, \mathbf{y}_{L0}) \leftarrow \phi_{\theta}(\mathbf{x})$
%     \STATE $\mathcal{L}_{\text{fine}} \leftarrow \mathrm{loss}(\mathbf{y}_{L0}, \mathbf{y})$
%     \STATE $\theta \leftarrow \theta - \eta \nabla_{\theta} \mathcal{L}_{\text{fine}}$
% \ENDFOR
% \end{algorithmic}
% \end{algorithm*}

% \begin{figure}[t]
%   \centering
%   \caption{Self-Conditioned Denoising Pretraining}
%   \label{fig:scd}
% \end{figure}



% ########## SINGLE COLUMN VERSION ######
\begin{algorithm}[t]
\small
% \setlength{\algorithmiclineskip}{2pt}
\renewcommand{\baselinestretch}{1.3}\selectfont
\caption{SCD Pretraining and Fine-tuning Pipeline}
\label{alg:pretrain_finetune}
\begin{algorithmic}[1]
\REQUIRE Dataset $\mathcal{D}$, noise scales $\sigma_{\text{corr}}, \sigma_{\text{reg}}$, model $\phi_{\theta}$
\ENSURE Trained parameters $\theta$

\vspace{0.5em}
\STATE \textbf{Pretraining}
\FOR{each minibatch $\mathbf{x} \sim \mathcal{D}$}
    \STATE $(\mathbf{x}, \tilde{\mathbf{x}}, \varepsilon) \leftarrow \textsc{Corrupt}(\mathbf{x}, \sigma_{\text{corr}}, \sigma_{\text{reg}})$ $\triangleright$ create views
    \STATE $(\mathbf{c}_{L0}, \mathbf{v}_{L1}, \mathbf{y}_{L0}) \leftarrow \phi_{\theta}(\mathbf{x})$ $\triangleright$ self-embedding
    \STATE $(\tilde{\mathbf{c}}_{L0}, \tilde{\mathbf{v}}_{L1}, \tilde{\mathbf{y}}_{L0}) \leftarrow \phi_{\theta}(\tilde{\mathbf{x}}, \mathbf{c}_{L0})$ $\triangleright$ conditional denoising
    \STATE $\mathcal{L}_{\text{pre}} \leftarrow \mathrm{MSE}(\tilde{\mathbf{v}}_{L1}, \varepsilon)$
    \STATE $\theta \leftarrow \theta - \eta \nabla_{\theta} \mathcal{L}_{\text{pre}}$
\ENDFOR

\vspace{0.5em}
\STATE \textbf{Fine-tuning / Inference}
\FOR{each labeled sample $(\mathbf{x}, \mathbf{y})$}
    \STATE $(\mathbf{c}_{L0}, \mathbf{v}_{L1}, \mathbf{y}_{L0}) \leftarrow \phi_{\theta}(\mathbf{x})$
    \STATE $\mathcal{L}_{\text{fine}} \leftarrow \mathrm{loss}(\mathbf{y}_{L0}, \mathbf{y})$
    \STATE $\theta \leftarrow \theta - \eta \nabla_{\theta} \mathcal{L}_{\text{fine}}$
\ENDFOR
\end{algorithmic}
\end{algorithm}

% #### DOUBLE COLUMN VERSION FIG #######
\begin{figure*}[t]
  \centering
   \includegraphics[width=0.85\linewidth]{scd_method_fig.png}
  \caption{Self-Conditioned Denoising Pretraining}
  \label{fig:scd}
\end{figure*}

% #### SINGLE COLUMN VERSION FIG #######
% \begin{figure}[t]
%   \centering
%   \caption{Self-Conditioned Denoising Pretraining}
%   \label{fig:scd}
% \end{figure}



\subsection{Backbone Architectures}
Following established baselines, we implement SCD using TorchMD-Net (also known as Equivariant Transformer or \textit{ET}) as the primary backbone \cite{TorchMD-Net_thölke2022, Frad_Ni2024, SliDe_Ni2025, UniCorn_Feng2024, NoisyNodes_Zaidi2023PVD}. To ensure a fair comparison, all models we train use an embedding dimension of 256, 8 layers, and 8 heads, consistent with previous SSL studies (see Table~\ref{tab:backbone_hparams}). 

The ET backbone architecture is a lightweight, attention-based GNN that lacks both tensor products and angular embeddings that are now commonly found in equivariant MLIPs. However, other denoising variations like \textit{Frad} \cite{Frad_Ni2024} and \textit{SliDe} \cite{SliDe_Ni2025} require explicit angular information and utilize a modified ``Geometric Equivariant Transformer'' (GET). To achieve a direct comparison with these methods, we evaluate SCD on both the standard ET and the more expressive GET backbones.

To incorporate conditional inputs into these non-conditional backbones, we replace existing pre-attention layer normalization with Adaptive Layer Normalization (\textit{AdaNorm}), analogous to its use in Diffusion Transformers \cite{DiT_peebles2023}. We find that AdaNorm has a negligible impact on inference speed or memory overhead. In contrast, the GET backbone is approximately $4\times$ slower and has twice as many parameters as the ET backbone. We denote the conditional versions of these two architectures as \textbf{CT} (Conditional TorchMD-Net) and \textbf{CGT} (Conditional Geometric TorchMD-Net), respectively. Additional implementation details and speed comparisons are provided in Appendix~\ref{Arch_and_implement_appendix}.

\subsection{Datasets}
Atomistic model training has often been siloed to single domains, typically focusing on small molecules, bio-molecules, or periodic materials independently. Although recent works have trained or pretrained on increasingly diverse datasets \cite{Atomica_Fang2025, ADIT_joshi2025, JMP_Shoghi2023}, multi-domain pretraining remains logistically challenging and under-explored. In this work, we evaluate models pretrained on a range of single domain datasets, with varying conformational diversity, as well as a combined multi-domain aggregate dataset (see Table~\ref{tab:dataset_summary}). All pretrained models are benchmarked on three representative tasks: 1) QM9 HOMO energy for small molecules \cite{QM9_Ramakrishnan2014}, 2) Matbench band gap for periodic materials \cite{matbench_2020}, and 3) Ligand Binding Affinity (LBA) for bio-molecules \cite{LBA_dataset}. Dataset descriptions are provided in Appendix~\ref{appendix_datasets}.


\begin{table*}[t]
    \centering
    \caption{Summary of Pretraining and Benchmarking Datasets}
    \label{tab:dataset_summary}
    \small
    \setlength{\tabcolsep}{4pt}
    \renewcommand{\arraystretch}{1.15}
    
    % \begin{adjustbox}{max width=\linewidth}
    \begin{tabular}{l l l l l}
    % \begin{tabularx}{\textwidth}{l X l l X}
   
    \toprule
         Pretraining Dataset&  Domain&  Unique Structures & Total Geometries& Origin/functional\\
         \midrule
         \textbf{PCQ} (PCQM4Mv2) & Small Molecules &  3.378M &  3.378M & B3LYP/6-31G* DFT\\
         \textbf{GEOM10} & Small Molecules &  304,339&  2.7M& GFN2-xTB , conformers\\
         \textbf{AMP20} (Alex-MP-20)  &  Periodic Materials &  607,673&  607,673& PBEsol DFT\\
         \textbf{SAIR} &  Protein-Ligand &    888k& 4.4M& Boltz-1x, conformers \\
         \textbf{ALL} (AllAtoms) &  Above, mixed & 5.2M & 11.3M & PCQ, GEOM10, AMP20, SAIR \\
         \textbf{OMol25} (4M split) &  mixed organic molecules&  ---  &  4M & MD, Non-equilibrium\\
         
        \midrule
        Benchmarking Dataset &  Domain &  Training Set Geometries &  \multicolumn{2}{l}{Task}\\
        \midrule
        \textbf{QM9} & Small Molecules  &  110k &  \multicolumn{2}{l}{properties computed from DFT}\\
        \textbf{Matbench} (mp gap) & Periodic Materials  & 85k &  \multicolumn{2}{l}{band gap energy computed from DFT} \\
        \textbf{LBA} & Protein-Ligand & 4k &  \multicolumn{2}{l}{ligand binding affinity ($-log(K_d) $) } \\
        % MD17 & Organic small molecules  &  10k &  \multicolumn{2}{l} {DFT Force/Energy}\\    
    \bottomrule
    \end{tabular}
    % \end{tabularx}
    % \end{adjustbox}
    
    % \label{dataset_summary}
    % \label{tab:placeholder}
    \footnotesize PCQ\cite{PubChemQC_Nakata2017}, GEOM\cite{GEOM_axelrod2022}, AMP20\cite{MaterialsProject, Alexandria_1, Alexandria_2}, SAIR \cite{Sair_dataset}, OMol25\cite{OMOL25_levine2025open},  Boltz-1x \cite{boltz1_2024}, QM9\cite{QM9_Ramakrishnan2014}, Matbench\cite{matbench_2020}, LBA\cite{LBA_dataset} 
\end{table*}

\section{Results}
\label{results}

% \subsection{Comparing SCD with existing SSL methods}
% \subsection{QM9: Small Molecule Property Prediction}
\subsection{SCD Outperforms Other SSL Methods on Atomistic Data}
Given the same pretraining dataset and backbone architecture, SCD pretraining significantly outperforms other forms of node denoising. When compared to standard node denoising pretraining (aka 'coord') we find that SCD conveys large improvements (19.6-45.5\%) across all QM9 targets tested, except heat capacity (Cv) (-5\%) (see table \ref{tab:scd_vs_coord}). In table \ref{tab:scd_vs_frad_slide} we compare SCD against SOTA methods Frad and Slide. We find that SCD pretraining, with either backbone architecture (CT or CGT) outperforms prior methods on most QM9 targets. Table \ref{tab:scd_vs_SSL_qm9} in the Appendix provides an additional comparison between our SCD pretrained models and other methods that used different backbone architectures and/or different pretraining datasets.


% ######### QM9: SCD vs Denoising ###############
\begin{table}[t]
\centering
\caption{\textbf{SCD vs. Standard Node Denoising}. Given the same backbone and pretraining datasets (PCQ), SCD significantly outperforms standard denoising (\textit{Coord}). Best results are \textbf{bolded}. All values reported as MAE on the test set.}
\label{tab:scd_vs_coord}
\small
\setlength{\tabcolsep}{5pt} % Adjusts padding for single-column fit
\begin{tabular*}{\columnwidth}{@{\extracolsep{\fill}}l r r r r}
\toprule
\textbf{Method} & \textbf{Baseline} & \textbf{Coord} & \textbf{SCD} & \textbf{\% imp.} \\
Backbone & ET & ET & CT & vs. Coord \\
\midrule
HOMO (meV) & 20.3 & 17.7 & \textbf{12.7} & 28.2\% \\
LUMO (meV) & 17.5 & 14.3 & \textbf{11.5} & 19.6\% \\
Gap (meV)  & 36.1 & 31.8 & \textbf{24.5} & 23.0\% \\
ZPVE (meV) & 1.84 & 1.71 & \textbf{1.18} & 21.0\% \\
$U_0$ (meV) & 6.15 & 6.57 & \textbf{3.58} & 45.5\% \\
$U$ (meV)   & 6.38 & 6.11 & \textbf{3.50} & 42.7\% \\
$H$ (meV)   & 6.16 & 6.45 & \textbf{3.52} & 45.4\% \\
$G$ (meV)   & 7.62 & 6.91 & \textbf{5.29} & 23.4\% \\
$\alpha$ ($a_0^3$) & 0.059 & 0.0517 & \textbf{0.0377} & 27.1\% \\
$C_v$ (cal/molK) & 0.026 & \textbf{0.020} & 0.021 & $-$5.0\% \\
\bottomrule
\end{tabular*}
\vspace{0.25em}
\end{table}
% ######### QM9: SCD vs Frad/slide SSL ###############
\begin{table}[t]
\centering
\caption{\textbf{SCD vs. Variations on Node Denoising}. All models are pretrained on PCQ. Percent improvement (\% imp.) is calculated between SliDe and SCD pretraining with the most similar backbone, CGT. \textbf{Bold} denotes best; \underline{underline} denotes second best. All values reported as MAE on the test set.}
\label{tab:scd_vs_frad_slide}
\small
\setlength{\tabcolsep}{3pt} % Tightens horizontal spacing
\begin{tabular*}{\columnwidth}{@{\extracolsep{\fill}}l r r r r r}
\toprule
\textbf{Method} & \textbf{Frad} & \textbf{SliDe} & \textbf{SCD} & \textbf{SCD} & \textbf{\% imp.}\\
Backbone & GET & GET & CT & CGT & vs. SliDe \\
\midrule
HOMO (meV) & 15.2 & 13.6 & \underline{12.7} & \textbf{9.65} & 29.0\% \\
LUMO (meV) & 13.7 & 12.3 & \underline{11.5} & \textbf{9.05} & 26.4\% \\
Gap (meV)  & 27.8 & 26.2 & \underline{24.5} & \textbf{19.7} & 24.8\% \\
ZPVE (meV) & 1.42 & 1.52 & \textbf{1.18} & \underline{1.22} & 19.8\% \\
$U_0$ (meV) & 5.33 & 4.28 & \textbf{3.58} & \underline{3.97} & 8.1\% \\
$U$ (meV)   & 5.62 & 4.29 & \textbf{3.50} & \underline{4.11} & 4.4\% \\
$H$ (meV)   & 5.55 & 4.26 & \textbf{3.52} & \underline{4.00} & 6.5\% \\
$G$ (meV)   & 6.19 & 5.37 & \textbf{5.29} & \underline{5.29} & 1.5\% \\
$\alpha$ ($a_0^3$) & \underline{0.0374} & \textbf{0.0366} & 0.0377 & 0.0383 & -4.6\% \\
$C_v$ (cal/molK) & 0.020 & \textbf{0.019} & 0.021 & \textbf{0.019} & 0.0\% \\
\bottomrule
\end{tabular*}
\vspace{0.25em}
\end{table}

\subsection{SCD Benefits From Pretraining on Any Atomistic Data Source}
% , particularly when applied to large, heterogeneous atomistic datasets
% Utilizing all the available atomistic data for pretraining has remained a challenge for foundation model development.
% Effective, low cost pretraining on large, diverse atomistic datasets that vary in quality or domain has been especially challenging for existing methods. 
Making use of all available atomistic data during pretraining remains a notable bottleneck in scaling the atomistic foundation models. Prior SSL work has been limited to ground-state geometries of small molecules\cite{Frad_Ni2024,SliDe_Ni2025, UniCorn_Feng2024}, while force-energy pretraining across datasets from different DFT functionals often requires the use of dataset specific heads or other specialize architecture components \cite{JMP_Shoghi2023,UMA_wood2025}. In this work, we find that SCD pretraining derived a significant benefit from every dataset tested. In Table \ref{tab:scd_on_conformers} we show how datasets from different domains, and conformational diversity, impact QM9 homo energy prediction error. Here, we report that SCD pretraining benefits from including both multiple conformers (GEOM1 vs GEOM10), and non-equilibrium structures (i.e. OMol25). We also find that large datasets are not necessarily required to see a significant benefit from SCD pretraining. Pretraining on a random 10\% subset of PCQ (300k samples) provided only slightly lower performance than pretraining on the whole dataset (3.4M samples). Overall, we find that the most significant predictor of downstream prediction accuracy is the similarity between the pretraining dataset and the finetuning dataset.
 
% ######### QM9: Training on a variaty of different datasets ###############
\begin{table*}[t]
\centering
\caption{\textbf{SCD Pretraining on Diverse Data Sources}. SCD enables effective pretraining across ground-state and non-equilibrium conformers from varied sources. Every pretraining configuration tested shows improvement over the non-pretrained baseline.}
\label{tab:scd_on_conformers}
\small
\begin{tabular*}{\textwidth}{@{\extracolsep{\fill}}l l r r r}
\toprule
\textbf{Pretraining Dataset} & \textbf{Source / Quality} & \textbf{Unique Mol.} & \textbf{Total Geom.} & \textbf{QM9 HOMO} \\ 
& & & & (MAE, meV) \\
\midrule
baseline &   &   &   & 20.3 \\ 
\midrule
\textbf{PCQ} & DFT (B3LYP/6-31G*) & 3,378,606 & 3,378,606 & \textbf{12.7} \\ 
\textbf{PCQ-300k} & DFT (B3LYP/6-31G*) & 300,000 & 300,000 & 13.0 \\ 
\textbf{GEOM10} & xTB (Conformers) & 304,339 & 2,791,929 & 13.2 \\ 
\textbf{GEOM1} & xTB (Lowest Energy) & 304,339 & 304,339 & 13.7  \\ 
\textbf{AMP20} & DFT (PBEsol) & 675,204 & 675,204 & 14.6 \\ 
\textbf{SAIR} & Boltz-1 (Synthetic Conformers) & 888,104 & 4.4M & 14.6 \\ 
% \textbf{SAIR-Ligand} & Boltz-2 (Conformers) & 888,104 & 4.4M & 14.9  \\ 
\textbf{SAIR-Pocket} & Boltz-1 (Synthetic Conformers)  & 888,104 & 4.4M & 14.5  \\ 
\textbf{ALL} & Mixed Aggregate & 5,246,253 & 11.3M & \textbf{12.7} \\ 
\textbf{OMol25}\textsuperscript{1} & Non-equilibrium (MD) & --- & 3,986,754 & 12.8 \\ 
\bottomrule
\end{tabular*}
\vspace{0.25em}
\flushleft\footnotesize \textsuperscript{1} The number of unique molecules in the OMol25 4M split is not reported \cite{OMOL25_levine2025open}.
\end{table*}

\subsection{Better Models Have Smoother Latent Spaces}
Figure \ref{fig:umap_mini} provides UMAPs for QM9 molecule embeddings from both untrained and pretrained models, colored by the number of atoms in each sample (N). We observe that an untrained ET model has a strong bias towards highly clustered, extensive representations - mean pooled atom embeddings that are highly correlated with the number of atoms in each molecule. We find that standard node denoising ('coord') results in a smooth but highly extensive representation space (this was also observed by \cite{GeomRep_li2024}). We believe this is a consequence of features that are too biased towards local node geometry, rather than global semantic representation (non-extensive features). In models pretrained by SCD, we observe greater smoothing and a reduced correlation with N, suggesting a more semantic representation. Empirically, we find models with a smoother, less extensive, embedding space reach a lower prediction error on QM9 targets. Additional plots are provided in figure \ref{fig:umap}. 


\begin{figure}[t]
  \centering
  \includegraphics[width=\linewidth]{umap_mini.png}
  % \includegraphics[width=0.85\linewidth]{umap_fig.png}
  \caption{\textbf{UMAP of molecule embeddings colored by atom count:} SCD results in smoother, less extensive, semantically rich representation. 'Coord' refers to standard node denoising pretraining. See figure \ref{fig:umap} for additional plots and details.
  % UMAP visualizations of embeddings for 10k QM9 molecules, colored by atom count ($N$). \textbf{(a)} Untrained GNNs exhibit a strong bias towards "extensive" representations, clustering molecules primarily by size ($N$). \textbf{(b)} Standard node denoising smooths the latent manifold but embeddings remain highly extensive. \textbf{(c--d)} SCD pretraining yields a smoother, less extensive (more semantic) embedding space. We observe that downstream performance improves as representations becomes smoother and less extensive (Performance: $a < b < c < d$). 'Molecule embeddings' are created by passing sum pooled atom embeddings through a two layer MLP 'embedding head'. In SCD 'Molecule embeddings' are used for self-embeddings. In our 'coord' model, the molecule 'embedding head' remains untrained.
  }
  \label{fig:umap_mini}
\end{figure}


\subsection{Comparing SCD Models with other Pretrained Models}
In table \ref{tab:scd_vs_jmp_qm9}, we show that pretraining and finetuning a simple, fast GNN with SCD is more cost effective than supervised training from scratch with more advanced tensor-product based architectures like Equiformer\cite{EquiformerV2_Liao2023}, yet it provides downstream results that are comparable to larger models pretrained by supervised force-energy prediction using orders of magnitude more data \cite{JMP_Shoghi2023}. 

SCD pretrained models also deliver strong results in property prediction for materials and proteins. In table \ref{tab:mpgap_brief} we show that our 10M parameter model pretrained on 675k periodic materials can achieve similar accuracy in bandgap prediction (within 2\%) as a 27M parameter model trained on 120M force-energy labeled geometries (108M of which are materials geometries). Table \ref{tab:lba_brief} presents a similar outcome for predictions of ligand binding affinity (LBA) on protein-ligand complexes. In this case, we can adjust the SCD objective to take advantage of an existing conditional relationship in the data: bound ligand conformation is determined by pocket geometry. We find that pretraining with 'pocket conditional ligand denoising' (CT-SCD-SAIR-Pocket) conveys an additional task specific benefit, and achieves state of the art results (See appendix section \ref{PairConditionalDenoising} for more details).
 
In order to make a direct comparison between SCD and force-energy pretraining, we also benchmark a force-energy pretrained CT model using the 4M split of OMol25, and compare it with a CT model pretrained using SCD on the exact same geometries (without force-energy labels). We summarize these results in Table \ref{results_summary}, and show that SCD pretraining performs roughly as well (-1.1\%) or slightly better (+3.9\%) than force-energy pretraining across all three tasks. 

% We summarize these results in Table \ref{results_summary}, and show that SCD pretrained models roughly match the [er] delivers similar performance to force-energy pretraining, outperforms 
% as well or better than force-energy pretraining across all three tasks. 

\begin{table*}[t]
\centering
\caption{\textbf{Comparing Computational Costs and Performance}. SCD pretraining enables competitive performance at a lower cost than more advanced Clebsch-Gordan tensor product based architectures, and large-scale force-energy (FE) pretrained models. \textbf{Bold} denotes best; \underline{underline} denotes second best.
% Notably, our models achieve performance comparable to large-scale JMP variants while using orders of magnitude fewer pretraining structures and GPU-hours. 
}
\label{tab:scd_vs_jmp_qm9}
\small
\begin{tabular*}{\textwidth}{@{\extracolsep{\fill}}l r r r r r r}
\toprule
\textbf{Metric} & \textbf{EquiformerV2} & \textbf{Baseline (ET)} & \textbf{CT-SCD} & \textbf{CGT-SCD} & \textbf{JMP-S} & \textbf{JMP-L} \\
% & \cite{EquiformerV2_Liao2023} & \cite{TorchMD-Net_thölke2022} & \textbf{(Ours)} & \textbf{(Ours)} & \cite{JMP_Shoghi2023} & \cite{JMP_Shoghi2023} \\
\midrule
Model Parameters & 11.2M & 7M & 10M & 17M & 27M & 235M \\
Pretraining Structures & 0 & 0 & 3.4M (U) & 3.4M (U) & 120M (L) & 120M (L) \\
Pretraining GPU-hours & 0 & 0 & \textbf{298} & \underline{1,308} & 5,700 & 34,400 \\
Finetuning GPU-h/task & 137 & --- & \textbf{46} & \underline{96} & --- & --- \\
Total hours (10 tasks) & \underline{1,370} & --- & \textbf{758} & 2,268 & $>5,700$ & $>34,400$ \\
\midrule
\textbf{QM9 Task (MAE)} & & & & & & \\
\midrule
HOMO (meV) & 14.0 & 20.3 & 12.7 & \underline{9.65} & 11.1 & \textbf{8.8} \\
LUMO (meV) & 13.0 & 17.5 & 11.5 & \underline{9.05} & 10.8 & \textbf{8.6} \\
Gap (meV) & 29.0 & 36.1 & 24.5 & \underline{19.7} & 23.1 & \textbf{19.1} \\
ZPVE (meV) & 1.47 & 1.84 & 1.18 & 1.22 & \underline{1.0} & \textbf{0.9} \\
$\alpha$ ($a_0^3$) & 0.050 & 0.059 & 0.0377 & 0.0383 & \underline{0.037} & \textbf{0.032} \\
$C_v$ (cal/mol K) & 0.023 & 0.026 & 0.021 & 0.019 & \underline{0.018} & \textbf{0.017} \\
$U_0$ (meV) & 6.17 & 6.15 & 3.58 & 3.96 & \underline{3.3} & \textbf{2.9} \\
$U$ (meV) & 6.49 & 6.38 & 3.50 & 4.11 & \underline{3.3} & \textbf{2.8} \\
$H$ (meV) & 6.22 & 6.16 & 3.52 & 4.00 & \underline{3.3} & \textbf{2.8} \\
$G$ (meV) & 7.57 & 7.62 & 5.29 & 5.29 & \underline{4.5} & \textbf{4.3} \\
\bottomrule
\end{tabular*}
\vspace{0.25em}
\flushleft\footnotesize (U) denotes unlabeled geometries; (L) denotes labeled force-energy structures.
\end{table*}
\begin{table}[h]
\centering
\caption{\textbf{Matbench Bandgap Prediction}. SCD models achieve competitive performance with significantly fewer parameters and unlabeled structures. Within each category best results are \textbf{bolded}; second best are \underline{underlined}.}
\label{tab:mpgap_brief}
\small
\setlength{\tabcolsep}{3pt}
\begin{tabular}{l c r r}
\toprule
\textbf{Model} & \textbf{MAE (eV)} & \textbf{Params} & \textbf{Data} \\
\midrule
\textbf{No Pretraining} \\
coGN \cite{CoGN_ruff2023} & \textbf{0.156} & --- & n/a \\ 
CT (Baseline) & \underline{0.186}\textsuperscript{\textdagger} & 10M & n/a \\
JMP-S \cite{JMP_Shoghi2023} & 0.235\textsuperscript{\textdagger} & 25.2M & n/a \\ 
JMP-L \cite{JMP_Shoghi2023} & 0.228\textsuperscript{\textdagger} & 235M & n/a \\ 
MODNet \cite{MODNet_DeBreuck2021} & 0.220 & --- & n/a \\ 
\midrule
\textbf{FE Pretrained} \\
JMP-L \cite{JMP_Shoghi2023} & \textbf{0.091} & 235M & 120M \\ 
JMP-S \cite{JMP_Shoghi2023} & \underline{0.121} & 27M & 120M \\
CT-FE-OMOL25 & 0.134\textsuperscript{\textdagger} & 10M & 4M \\
HackNIP \cite{HackNIP_kim2025} & 0.150 & $>25$M & $>32$M\textsuperscript{*} \\
\midrule
\textbf{SCD Pretrained (Ours)} \\
CT-SCD-AMP20 & \textbf{0.123} & 10M & 675k \\
CT-SCD-ALL & \underline{0.132}\textsuperscript{\textdagger} & 10M & 11.3M \\
CT-SCD-OMOL25 & 0.136\textsuperscript{\textdagger} & 10M & 4M \\
\bottomrule
\end{tabular}
\vspace{0.25em}
\flushleft\footnotesize \textsuperscript{\textdagger} Reported from Fold 0. See full results in table \ref{tab:mpgap_full}. \\
\textsuperscript{*} Estimated from Orb-v2 backbone.
\end{table}
% % ########################### LBA Single column summary table ####################################
% \begin{table}[h]
% \centering
% \caption{\textbf{Ligand binding affinity:} Results (RMSE) from Models evaluated on the 30\% sequence identity split. SCD models achieve parity with supervised SOTA despite significantly smaller parameter counts. Best results are \textbf{bolded}; second best are \underline{underlined}. }
% \label{tab:lba_brief}
% \small
% \begin{tabular*}{\columnwidth}{@{\extracolsep{\fill}}l r r}
% \toprule
% \textbf{Model} & \textbf{RMSE} $\downarrow$ & \textbf{Params} \\
% \midrule
% \textbf{No Pretraining} \\
% % ProtNet-All-Atom \cite{ProtNet_wang2023} & 1.463 & --- \\
% ProtNet \cite{ProtNet_wang2023} & 1.463 & --- \\
% EPT-Scratch \cite{EPT_jiao2025} & \textbf{1.378} & 30M \\
% CT (Baseline) & 1.510 & 10M \\
% \midrule
% \textbf{Other Pretrained} \\
% Uni-Mol \cite{UniMol_Zhou2023} & 1.520 & 47.6M \\
% EPT-Protein \cite{EPT_jiao2025} & 1.329 & 30M \\
% EPT-MultiDomain \cite{EPT_jiao2025} & \underline{1.322} & 30M \\
% ADiT-S \cite{lba-adit} & 1.337 & 12M \\
% ADiT-M \cite{lba-adit} & 1.353 & 35M \\
% ADiT-L \cite{lba-adit} & \textbf{1.308} & 253M \\
% \midrule
% \textbf{SCD Pretrained (Ours)} \\
% CT-SCD-ALL & 1.372 & 10M \\
% CT-SCD-SAIR & \underline{1.337} & 10M \\
% CT-SCD-SAIR-Pocket & \textbf{1.304} & 10M \\
% CT-SCD-OMOL25 & 1.389 & 10M \\
% \bottomrule
% \end{tabular*}
% \footnotesize Full results provided in Table \ref{tab:lba_full}.
% \end{table}


% % ########################### LBA Single column summary table ####################################
\begin{table}[h]
\centering
\caption{\textbf{Ligand binding affinity:} RMSE on the 30\% (id30) and 60\% (id60) sequence identity overlap splits. Best results are \textbf{bolded}; second best are \underline{underlined}. }
\label{tab:lba_brief}
\small
\begin{tabular*}{\columnwidth}{@{\extracolsep{\fill}}l r r}
\toprule
\textbf{Model} & \textbf{RMSE: id30 / id60} $\downarrow$ & \textbf{Params} \\
\midrule
\textbf{No Pretraining} \\
% ProtNet-All-Atom \cite{ProtNet_wang2023} & 1.463 & --- \\
ProtNet \cite{ProtNet_wang2023} & 1.463 / 1.343 & --- \\
EPT-Scratch \cite{EPT_jiao2025} & 1.378 / 1.277 & 30M \\
CT (Baseline) & 1.510 / 1.386 & 10M \\
\midrule
\textbf{Other Pretrained} \\
Uni-Mol \cite{UniMol_Zhou2023} & 1.520 / 1.619 & 47.6M \\
EPT-Protein \cite{EPT_jiao2025} & 1.329 / 1.235 & 30M \\
EPT-Multi \cite{EPT_jiao2025} & 1.322 / 1.227 & 30M \\
ADiT-S \cite{lba-adit} & 1.337 / 1.413 & 12M \\
ADiT-M \cite{lba-adit} & 1.353 / 1.335 & 35M \\
ADiT-L \cite{lba-adit} & \underline{1.308} / 1.246 & 253M \\
\midrule
\textbf{SCD Pretrained (Ours)} \\
CT-SCD-ALL & 1.372 / 1.218 & 10M \\
CT-SCD-PCQ & 1.332 / 1.226 & 10M \\
CT-SCD-SAIR & 1.337 / \underline{1.196} & 10M \\
CT-SCD-SAIR-Pocket & \textbf{1.304} / 1.200 & 10M \\
CT-SCD-OMOL25 & 1.389 / \textbf{1.175} & 10M \\
\bottomrule
\end{tabular*}
\footnotesize Full results provided in Table \ref{tab:lba_full}.
\end{table}



% ############ SUMMARY TABLE ############
\begin{table*}[t]
\centering
\caption{\textbf{Comparing the Impact of Different Pretraining Data}. Pretraining on any data source benefits every task, though downstream performance gains are greatest when pretraining on in-domain data. Datasets with high elemental diversity yield the highest average improvements. \textit{CT-FE} models denote supervised force-energy pretraining; \textit{CT-SCD} denotes self-supervised pretraining. Best results are \textbf{bold}; second best are \underline{underlined}.}
\label{results_summary}
\small
\begin{tabular*}{\textwidth}{@{\extracolsep{\fill}}l r r r r r r r}
\toprule
 & \multicolumn{2}{c}{\textbf{QM9 (Small Mol.)}} & \multicolumn{2}{c}{\textbf{Matbench (Crystals)}} & \multicolumn{2}{c}{\textbf{LBA (Proteins)}} & \\
\cmidrule(lr){2-3} \cmidrule(lr){4-5} \cmidrule(lr){6-7}
\textbf{Model} & \textbf{HOMO} & \textbf{\% imp.} & \textbf{MP Gap} & \textbf{\% imp.} & \textbf{RMSE} & \textbf{\% imp.} & \textbf{Avg \%} \\
(Pretraining Data) & (meV) & & (eV) & & (id60) & & \textbf{imp.} \\
\midrule
ET/CT (Baseline) & 20.3 & 0.0\% & 0.186 & 0.0\% & 1.386 & 0.0\% & 0.0\% \\
\midrule
CT-SCD-PCQ & \textbf{12.7} & 37.4\% & 0.174 & 6.5\% & 1.226 & 11.5\% & 18.5\% \\
CT-SCD-GEOM10 & \underline{13.2} & 35.0\% & 0.177 & 4.8\% & 1.211 & 12.6\% & 17.5\% \\
CT-SCD-AMP20 & 14.6 & 28.1\% & \textbf{0.122} & 34.4\% & 1.219 & 12.0\% & 24.8\% \\
CT-SCD-SAIR & 14.6 & 28.1\% & 0.182 & 2.2\% & 1.196 & 13.7\% & 14.6\% \\
CT-SCD-ALL & \textbf{12.7} & 37.4\% & \underline{0.132} & 29.0\% & 1.218 & 12.1\% & \underline{26.2}\% \\
\midrule
CT-FE-OMol25 & 13.6 & 33.0\% & 0.134 & 28.0\% & \underline{1.187} & 14.4\% & 25.1\% \\
CT-SCD-OMol25 & 12.8 & 36.9\% & 0.136 & 26.9\% & \textbf{1.175} & 15.2\% & \textbf{26.4}\% \\
\bottomrule
\end{tabular*}
\end{table*}


\subsection{Conclusion}

In this work, we present a limited exploration of the SCD pretraining objective across multiple domains of pretraining data and property prediction tasks.  We find that SCD pretraining provides robust improvements on all tasks investigated and matches or outperforms the benefits of force-energy pretraining, as well as other pretraining methods. Small SCD pretrained models achieve state-of-the-art results across multiple domains: QM9,  Matbench properties, and Ligand binding affinity benchmarks. Overall, our findings can be summarized through several key observations:

\textbf{All data is useful for pretraining:} We observed that SCD pretraining on atomistic data from any domain or source (DFT accurate ground-state, 'non-equilibrium', or diffusion-generated) conveyed an improvement across all benchmarking tasks. Surprisingly, we find that even ligand binding affinity predictions are improved when pretraining on periodic materials (+12\%), and pretraining with non-periodic, non-equilibrium molecule geometries (OMol25) conveyed a surprising benefit to bandgap predictions in periodic materials (+26.9\%) (see Table \ref{results_summary}). However, we hypothesize that at least some of these benefits may stem from pre-smoothing of the representation space prior to task-specific fine-tuning (Figure \ref{fig:umap}).

\textbf{Pretraining is not necessarily expensive:}
Our backbone architecture is neither complex nor state-of-the-art, and without pretraining, it provides only a modest baseline. However, we find that a simple, fast GNN pretrained by SCD can outperform supervised learning with advanced tensor-product based architectures (e.g., EquiformerV2 \cite{EquiformerV2_Liao2023}) in both prediction accuracy and total compute cost (including pretraining GPU-hours. Table \ref{tab:scd_vs_jmp_qm9}).  Furthermore, large datasets are not always necessary to realize these benefits. Pretraining on a random 10\% slice of the PCQ dataset (300k molecules) captures 97\% of the performance gain of the full dataset for QM9 tasks, and pretraining on just 607k materials yields band gap accuracies within 1.7\% of larger models pretrained on 120M force-energy labeled geometries, 108M of which are from materials. Overall, we find that efficient pretraining is more impactful and cost effective than implementing inductive bias through complex architectural elements (as has been found in other fields). 

\textbf{Data relevance is key to performance:}
Across all tasks explored, the strongest predictor of success was similarity between samples in the pretraining dataset and the fine-tuning dataset. This was most apparent in materials bandgap prediction, in which pretraining datasets containing inorganic atoms provided far greater benefit than those containing only organic elements. Notably, our model pretrained on a combine super-set of single domain data (CT-SCD-ALL) achieved significantly better average performance across all domains, but did not surpass the performance of copies pretrained on single-domain data for in-domain tasks. We attribute this small drop in performance to the limited expressivity of our 10M parameter model rather than representation conflicts in multi-domain training, though further studies are needed to confirm this. We expect larger models to receive greater benefit from pretraining with large multi-domain datasets.


\textbf{Pretraining a Foundation Model does not require labels:}
Force-energy pretraining has proven to be an effective and scalable pretraining strategy when DFT labels are readily available \cite{JMP_Shoghi2023, HackNIP_kim2025}. However, computing large numbers of new ground truth (DFT) forces and energies is extremely expensive and time-consuming. Finding effective methods of pretraining with unlabeled, low fidelity, and or cheaply generated atomistic geometries could enable further generalization of atomistic foundation models without requiring a significant expansion of DFT labeled datasets. We provide SCD as a demonstration that this is possible through self-supervised learning. When controlling for architecture and pretraining geometries (labeled vs unlabeled), SCD pretraining can match or surpass the performance of supervised force-energy pretraining in downstream property prediction tasks. Further, we show that even diffusion generated conformers (i.e. SAIR) benefit SCD pretraining. As generative methods in the atomistic sciences continue to advance, we anticipate an increased availability of large 'synthetic' datasets for pretraining. 

In developing SCD, our aim is to inspire further exploration into self-supervised learning and efficient pretraining for atomistic foundation models. More work is still needed to investigate the effects of scaling across dataset and model dimensions. To assist in this effort, we've made our code and pre-trained models publicly available at \url{https://github.com/TyJPerez/SelfConditionedDenoisingAtoms}

% intend to make our code, including dataset loaders, training algorithms, and pretrained model weights, available online.
% % TODO Add links to github, etc

% Our code and pretrained models are publicly available at \url{https://github.com/TyJPerez/SelfConditionedDenoisingAtoms}





% #######  FOR PREPRINT ONLY
% Acknowledgements should only appear in the accepted version.
\section*{Acknowledgements}
We thank Kaiming He and Bowen Deng for valuable feedback and helpful discussions on
this work. This research used resources of the National Energy Research Scientific
Computing Center (NERSC), a Department of Energy User Facility using NERSC award
ALCC-ERCAP 0038200.

% #######  FOR PREPRINT ONLY

\section*{Impact Statement}
This paper aims to advance applications of Machine Learning in the atomistic sciences. Specifically, we focus on improving the efficacy of self-supervised pretraining, potentially reducing the need for computationally expensive DFT labels. Any societal consequences or impacts that typically relate to the development of atomistic foundation models also apply here. Foundation models and their derivative tools could lower the cost of discovery in many areas of science, including dual use technologies. 

% This paper presents work intended to advance 
% whose goal is to advance the field of Machine
% Learning. There are many potential societal consequences of our work, none
% which we feel must be specifically highlighted here.

% Authors are \textbf{required} to include a statement of the potential broader
% impact of their work, including its ethical aspects and future societal
% consequences. This statement should be in an unnumbered section at the end of
% the paper (co-located with Acknowledgements -- the two may appear in either
% order, but both must be before References), and does not count toward the paper
% page limit. In many cases, where the ethical impacts and expected societal
% implications are those that are well established when advancing the field of
% Machine Learning, substantial discussion is not required, and a simple
% statement such as the following will suffice:

% ``This paper presents work whose goal is to advance the field of Machine
% Learning. There are many potential societal consequences of our work, none
% which we feel must be specifically highlighted here.''

% The above statement can be used verbatim in such cases, but we encourage
% authors to think about whether there is content which does warrant further
% discussion, as this statement will be apparent if the paper is later flagged
% for ethics review.

\bibliography{example_paper}
\bibliographystyle{icml2026}

%%%%%%%%%%%%%%%%%%%%%%%%%%%%%%%%%%%%%%%%%%%%%%%%%%%%%%%%%%%%%%%%%%%%%%%%%%%%%%%
%%%%%%%%%%%%%%%%%%%%%%%%%%%%%%%%%%%%%%%%%%%%%%%%%%%%%%%%%%%%%%%%%%%%%%%%%%%%%%%
% APPENDIX
%%%%%%%%%%%%%%%%%%%%%%%%%%%%%%%%%%%%%%%%%%%%%%%%%%%%%%%%%%%%%%%%%%%%%%%%%%%%%%%
%%%%%%%%%%%%%%%%%%%%%%%%%%%%%%%%%%%%%%%%%%%%%%%%%%%%%%%%%%%%%%%%%%%%%%%%%%%%%%%
\newpage
\appendix
\onecolumn

% \section{You \emph{can} have an appendix here.}
% You can have as much text here as you want. The main body must be at most $8$
% pages long. For the final version, one more page can be added. If you want, you
% can use an appendix like this one.

% The $\mathtt{\backslash onecolumn}$ command above can be kept in place if you
% prefer a one-column appendix, or can be removed if you prefer a two-column
% appendix.  Apart from this possible change, the style (font size, spacing,
% margins, page numbering, etc.) should be kept the same as the main body.




% \subsection{Why Does Self-Supervised Learning Matter for 3D Atomistic data?}
% In other fields, the paradigm shift towards foundation models would not have been possible without the development of self-supervised learning  objectives that enabled pretraining on vast, heterogeneous, and largely unlabeled datasets. However, the complex structure of 3D atomistic data poses a unique challenge for the development of effective self-supervised learning objectives. Further, in physical chemistry and materials science, 3D geometries are often generated using expensive DFT relaxations or simulations, which also yield corresponding force and energy labels. Consequently, the fraction of 3D geometries with force–energy labels is relatively high compared to other data modalities such as text or images. This abundance has made large-scale supervised pretraining a feasible path for progress in the field, driving consistent improvements with scale. However, the relative prevalence of "ground-truth" (DFT based) force-energy labels is unlikely to remain the case in the near future as emerging foundation potentials and diffusion-based generative models provide far more efficient routes to obtaining 3D geometries than DFT.

\section{Practical Notes on Node Denoising }
\label{denoising_details}

\subsection{Interpreting Noise Scales} 
We find that the scale of noise applied to atomistic structures plays a significant role in its utility as a pretraining or regularization tool. Here we formalize two distinct regimes:

\textbf{Regularizing noise} utilizes a small scale with respect to bond length (single bonds range from $\sim$ 1.0 to 2.0\AA ), such as $\sigma \approx 0.005\text{\AA}$), in order to reduce overfitting. At this scale gaussian noise does not alter single-double bond distinguishability, and is often used as an auxiliary task \cite{NoisyNodes_Godwin2022,Equiformer1_Liao2022, EquiformerV2_Liao2023}. We find that regularizing noise improves downstream results when used in SCD pretraining (on non-corrupted inputs) and performance during some finetuning tasks, but not all.

\textbf{Corruption noise} used for pretraining, is typically one order of magnitude larger ($\sigma \approx 0.04\text{\AA}$), and often perturbs molecular geometries enough to disrupt single/double/triple bond identities (+/- 0.1\AA), but without moving atoms far enough to be considered bond breaking. We find that noise scales much larger than this harm pretraining performance. 


\subsection{Reinterpreting Denoising for Atomistic Systems}
Many previous works have suggested a relationship between denoising objectives and learning `force fields'; however, we do not find this to be a useful hypothesis for why denoising benefits representation. It is often implied that the `field' learned by noise prediction is beneficial due to its similarity to DFT forces, yet prior work has shown that noise vectors on ground-state geometries correlate poorly with true DFT forces \cite{SliDe_Ni2025}. Consistent with this, we find that for our models trained on ground-truth DFT forces, adding an auxiliary node denoising objective harms performance on property prediction, suggesting a low similarity between these targets (see Table \ref{tab:fe_finetune}).

Instead, we hypothesize that denoising objectives succeed because they encourage the development of smooth, distinguishable local atom embeddings. SCD complements this by introducing a global semantic objective. Through this lens, we attribute the transferability of force-energy pretraining to its demand for high geometric sensitivity and representational complexity, rather than to its explicit connection to physical principles. Force prediction encourages equivariant sensitivity to local perturbations, while energy prediction promotes global representation. Based on this hypothesis, we suggest that self-supervised learning methods should employ reconstruction or de-corruption objectives that challenge models with both local and global objectives similar to force-energy pretraining and SCD. 



\section{Backbone Architecture and Implementation}
\label{Arch_and_implement_appendix}

Self-Conditioned Denoising pretraining can be implemented with any architecture designed for diffusion. To modify an existing non-conditional GNN for SCD pretraining, we replace the pre-norm of each message passing layer with an adaptive layer norm (AdaNorm), similar to a diffusion transformer \cite{DiT_peebles2023}. 

figure \ref{fig:adanorm} diagrams our implementation of AdaNorm used for SCD. Because SCD pretraining requires two forward passes per-step we find that including drop path with a drop probability of 0.1 improves training stability and provides an additional form of regularization. Our drop path is implemented to jointly drop both L0 and L1 embeddings. During pretraining we also freeze atom type/element embeddings. Without freezing element embeddings, we find that element embeddings become vanishingly small, leading to downstream instability. On github \cite{UnstableCoord_Zaidi2023GitHubIssue}, previous authors have noted challenges (gradient explosion) during the downstream fine tuning of models trained by standard node-denoising pretraining (aka 'coord') - likely a result of small or zeroed out embedding vectors.


Table ~\ref{tab:arch_speed_mem} shows a breakdown of parameter count, speed, and memory usage for the architectures used in this work. Each architecture shown uses the same embedding size (256), number of layers (8), and number of attention heads (8). CT and CGT correspond to the modified architectures of ET and GET that incorporate conditional layer norms. In this table, we can see that adding conditional guidance to an architecture has little effect on its speed and memory usage as compared to the addition of edge updates and angle embeddings. CT is overall 4 times faster, almost half the size, and has half the peak memory usage as GET. 

% \begin{table}[t]
% \centering
% \caption{Architecture speed and memory comparison on QM9 samples (batch size 128, averaged over 32 batches on an RTX~A5000). All models listed have the same embedding dimension, number of layers, and number of attention heads}
% \label{tab:arch_speed_mem}

% % --------- Subtable A: raw metrics ----------
% \begin{subtable}{\linewidth}
% \centering
% \caption{Latency, throughput, and memory (absolute).}
% \label{tab:arch_speed_mem_abs}
% \begin{adjustbox}{max width=\linewidth}
% \begin{tabular}{l c c c c c}
% \toprule
% \textbf{Arch} &
% \textbf{Backbone Params (M)} &
% \shortstack[c]{\textbf{forward pass (ms)}\\\textbf{(Eval / Train)}} &
% \shortstack[c]{\textbf{samples/sec}\\\textbf{(Eval / Train)}} &
% \shortstack[c]{\textbf{Avg Mem (GB)}\\\textbf{(Eval / Train)}} &
% \shortstack[c]{\textbf{Peak Mem (GB)}\\\textbf{(Eval / Train)}} \\
% \midrule
% ET   & 6.53  & 65.7 / 69.13   & 1946.4 / 879.7 & 0.040 / 0.128 & 1.07 / 8.47 \\
% CT  & 9.17  & 70.59 / 74.5   & 1811.6 / 802.1 & 0.050 / 0.148 & 1.08 / 9.00 \\GET (Frad/SliDe) & 13.39 & 264.9 / 369.5 & 483.1 / 193.8 & 0.067 / 0.18  & 1.62 / 20.37 \\
% CGT & 16.03 & 262.8 / 379.3 & 486.8 / 189.0 & 0.076 / 0.20  & 1.63 / 20.72 \\
% \bottomrule
% \end{tabular}
% \end{adjustbox}
% \end{subtable}

% \vspace{0.6em}
% % --------- Subtable B: relative metrics ----------
% \begin{subtable}{\linewidth}
% \centering
% \caption{Relative metrics (normalized to ET).}
% \label{tab:arch_speed_mem_rel}
% \begin{adjustbox}{max width=\linewidth}
% \begin{tabular}{l c c c}
% \toprule
% \textbf{Arch} &
% \shortstack[c]{\textbf{Rel. Mem to ET}\\\textbf{(Eval / Train)}} &
% \shortstack[c]{\textbf{Rel. Speed to ET}\\\textbf{(Eval / Train)}} &
% \textbf{Rel. Size to ET} \\
% \midrule
% ET   & 1.0x / 1.0x  & 1.0x / 1.0x  & 1.0x  \\
% CT  & 1.25x / 1.15x & 0.92x / 0.93x & 1.4x  \\
% GET (Frad/SliDe) & 1.67x / 1.38x & 0.25x / 0.19x & 2.0x  \\
% CGT & 1.90x / 1.53x & 0.25x / 0.18x & 2.45x \\
% \bottomrule
% \end{tabular}
% \end{adjustbox}
% \end{subtable}
% \vspace{0.25em}
% \end{table}


\begin{table}[t]
\centering
\caption{\textbf{Architecture Efficiency Comparison}. Benchmarked using QM9 molecules, with batch size 128, averaged over 32 batches, on a NVIDIA RTX A5000 GPU. All models share identical embedding dimensions, layer counts, and attention heads. Parameter counts in this table encompass only backbone architecture, excluding the embedding and prediction heads. All metrics shown are computed from singular forward passes.}
\label{tab:arch_speed_mem}
\small

% --------- Subtable A: Absolute Metrics ----------
\begin{subtable}{\linewidth}
\centering
\caption{\textbf{Absolute Metrics}: Latency (lower is better), Throughput (higher is better), and Memory.}
\label{tab:arch_speed_mem_abs}
\setlength{\tabcolsep}{4pt}
\begin{tabular}{l c cc cc cc cc}
\toprule
\textbf{Model} & \textbf{Params} & \multicolumn{2}{c}{\textbf{Latency (ms)} $\downarrow$} & \multicolumn{2}{c}{\textbf{Throughput (samp/s)} $\uparrow$} & \multicolumn{2}{c}{\textbf{Avg Mem (GB)}} & \multicolumn{2}{c}{\textbf{Peak Mem (GB)}} \\
\cmidrule(lr){3-4} \cmidrule(lr){5-6} \cmidrule(lr){7-8} \cmidrule(lr){9-10}
& (M) & Eval & Train & Eval & Train & Eval & Train & Eval & Train \\
\midrule
ET & 6.5 & 65.7 & 69.1 & 1946 & 880 & 0.04 & 0.13 & 1.07 & 8.47 \\
CT & 9.2 & 70.6 & 74.5 & 1812 & 802 & 0.05 & 0.15 & 1.08 & 9.00 \\
GET & 13.4 & 264.9 & 369.5 & 483 & 194 & 0.07 & 0.18 & 1.62 & 20.37 \\
CGT & 16.0 & 262.8 & 379.3 & 487 & 189 & 0.08 & 0.20 & 1.63 & 20.72 \\
\bottomrule
\end{tabular}
\end{subtable}

\vspace{1em}

% --------- Subtable B: Relative Metrics ----------
\begin{subtable}{\linewidth}
\centering
\caption{\textbf{Relative Metrics}: Normalized to ET baseline ($1.0\times$).}
\label{tab:arch_speed_mem_rel}
\setlength{\tabcolsep}{8pt}
\begin{tabular}{l c cc cc}
\toprule
\textbf{Model} & \textbf{Param Ratio} & \multicolumn{2}{c}{\textbf{Rel. Speed} $\uparrow$} & \multicolumn{2}{c}{\textbf{Rel. Avg Mem} $\downarrow$} \\
\cmidrule(lr){3-4} \cmidrule(lr){5-6}
& & Eval & Train & Eval & Train \\
\midrule
ET & $1.00\times$ & $1.00\times$ & $1.00\times$ & $1.00\times$ & $1.00\times$ \\
CT & $1.40\times$ & $0.92\times$ & $0.93\times$ & $1.25\times$ & $1.15\times$ \\
GET & $2.00\times$ & $0.25\times$ & $0.19\times$ & $1.67\times$ & $1.38\times$ \\
CGT & $2.45\times$ & $0.25\times$ & $0.18\times$ & $1.90\times$ & $1.53\times$ \\
\bottomrule
\end{tabular}
\end{subtable}
\end{table}

\begin{figure}[t]
  \centering
  \includegraphics[width=0.6\linewidth]{arch_both_res.png}
  \caption{\textbf{Adaptive layer normalization for SCD:} On the left we illustrate a simplified diagram of the original residual layer block in TorchMD-Net (ET). Here, 'X' represents invariant (L0) embeddings while 'V' represents equivariant vector embeddings (L1), and 'C' represents the conditioning embedding. On the right, we show how we modify the original residual block with a conditional scale, shift, and gate of the invariant embeddings. In the gate block the gamma vector is attenuated by a tanh function and multiplied to the output of the ET Attention block. For this work, we use a two layer MLP for conditioning in each layer.}
  \label{fig:adanorm}
\end{figure}


\begin{table}[t]
\centering
\caption{Backbone architecture hyperparameters used in this work. CT-small is only used in the MD17 benchmark}
\label{tab:backbone_hparams}
\begin{tabular}{lccc}
\toprule
\textbf{Hyperparameter} & \textbf{CT} & \textbf{CGT} & \textbf{CT-small} \\
\midrule
Layers                & 8   & 8   & 6   \\
Attention heads       & 8   & 8   & 8   \\
Embedding dimension   & 256 & 256 & 128 \\
Layer pooling         & sum & sum & sum \\
Graph cutoff (\AA)    & 5   & 5   & 5   \\
L0 head pooling       & sum & sum & sum \\
\bottomrule
\end{tabular}
\end{table}

% \FloatBarrier % ############################

\section{Supplementary Results}


% ########### UMPA figure
% \begin{figure*}[t]
%   \centering
%   % Make sure the path to your image is correct relative to your main tex file
%   % \caption{\textbf{UMAP visualization of pretrained model embeddings.} Embeddings for 10k QM9 molecules are colored by atom count ($N$). Untrained GNNs show strong clustering biased by simple features like $N$. Models with higher downstream performance exhibit smoother embedding manifolds with less trivial correlation to $N$. Downstream performance improves from left to right: (a) $<$ (b) $<$ (c) $<$ (d).}
%   % \label{fig:umap_wide}
  
%   \includegraphics[width=\linewidth]{umap_fig.png}
%   \caption{UMAP of untrained and pretrained model embeddings for 10k QM9 molecules. Points are colored by number of atoms (N) per molecule. Row a) shows that untrained GNNs are biased towards over clustering molecular embeddings at high correlation with number of nodes/atoms (N) - 'extensive representation'. In rows b, node denoising smooths the representation space, but does not adequately address the bias towards extensive representation. Rows c and d show that embeddings from SCD model are less extensive, and smoother. We find that pretrained models with smoother, less extensive representation provide better downstream performance in property prediction. Performance trend: worst $a<b<c<d$ best.}
%   \label{fig:umap}
% \end{figure*}

% \begin{figure*}[t]
%   \centering 
%   \includegraphics[width=\linewidth]{umap_fig.png}
%   \caption{UMAP of untrained and pretrained model embeddings for 10k QM9 molecules. Points are colored by number of atoms (N) per molecule. Row a) shows that untrained GNNs are biased towards over clustering molecular embeddings at high correlation with number of nodes/atoms (N) - 'extensive representation'. In rows b, node denoising smooths the representation space, but does not adequately address the bias towards extensive representation. Rows c and d show that embeddings from SCD model are less extensive, and smoother. We find that pretrained models with smoother, less extensive representation provide better downstream performance in property prediction. Performance trend: worst $a<b<c<d$ best.}
%   \label{fig:umap}
% \end{figure*}

\begin{figure*}[t]
  \centering
  % \includegraphics[width=\linewidth]{umap_fig.png}
  \includegraphics[width=0.85\linewidth]{umap_fig.png}
  \caption{\textbf{Pretraining Effects on the Embedding Space.} UMAP visualizations of embeddings for 10k QM9 molecules, colored by atom count ($N$). \textbf{(a)} Untrained GNNs exhibit a strong bias towards "extensive" representations, clustering molecules primarily by size ($N$). \textbf{(b)} Standard node denoising smooths the latent manifold but embeddings remain highly extensive. \textbf{(c--d)} SCD pretraining yields a smoother, less extensive (more semantic) embedding space. We observe that downstream performance improves as representations becomes smoother and less extensive (Performance: $a < b < c < d$). 'Molecule embeddings' are created by passing sum pooled atom embeddings through a two-layer MLP 'embedding head'. In the case of 'ET-Coord' the embedding head is untrained. In SCD 'Molecule embeddings' are used for self-embeddings. In our 'coord' model, the molecule 'embedding head' remains untrained.}
  \label{fig:umap}
\end{figure*}



\subsection{Ligand Binding Affinity}
\subsubsection{Conditional Denoising of Coupled Components: Pockets and Ligands}
\label{PairConditionalDenoising}
When dataset samples ($\mathbf{x}$) contain interdependent pairs ($\mathbf{x}=A\cup B$), such as a protein pocket and a bound ligand, we can alter the self-conditioned denoising objective to instead denoise one component based on an embedding from the other (eqn \ref{component_conditional_denoising}). In doing so we can exploit known conditional relationships in the dataset during pretraining.

\begin{equation}\label{component_conditional_denoising}
\begin{aligned}
&\mathbb{E}_{q_{\sigma}(\tilde{\mathbf{x}}, \mathbf{x})}
\left[
  \left\lVert
    \phi_{\theta}(\tilde{\mathbf{x}}|\mathbf{c}) - \varepsilon
  \right\rVert_{2}^{2}
\right]
= \mathbb{E}_{q_{\sigma}(\tilde{\mathbf{A}}, \mathbf{B})}
\left[
  \left\lVert
    \phi_{\theta}(\tilde{\mathbf{A}}|\phi_{\theta}(\mathbf{B})_{L0})_{L1} - \varepsilon_A
  \right\rVert_{2}^{2}
\right]\\
% &\text{where } \mathbf{x} = \{\mathbf{A}, \mathbf{B}\} \\
\end{aligned}
\end{equation}


In table \ref{tab:lba_full} we present results from three variations of pretraining on SAIR. In CT-SCD-SAIR, we pretrain the model using canonical SCD, embedding and denoising the entire pocket-ligand complexes together. However, in pretraining CT-SCD-SAIR-Lig the ligand structure is embedded and used to conditionally denoise the pocket atoms; and in CT-SCD-SAIR-Pocket we do the reverse, pocket atoms are embedded and used to conditionally denoise the ligand atoms. In doing this we find that CT-SCD-SAIR-Pocket performs the best (specifically on low sequence overlap), CT-SCD-SAIR the second best, and CT-SCD-SAIR-Lig the worst. This follows biochemical intuition, as we'd expect ligand conformation to be strongly dependent on pocket geometry, but not the other way around. We hold this as an example of how exploiting known conditional relationships can benefit learned representation during pretraining. We suspect that component conditional de-corrupt objectives may also benefit other pairwise tasks.
% ### larger results table
% ######### LBA: FULL Results ############### --- Appendix
\begin{table*}[t]
\centering
\caption{\textbf{SCD on Proteins}: Ligand Binding Affinity (LBA) predictions over two sequence identity splits (30\% and 60\%). SCD models pretrained on synthetic or diverse data (e.g., SAIR, OMol25) can outperform larger models pretrained on PDB-based datasets. Best results within each category are \textbf{bolded}; second best are \underline{underlined}.}
\label{tab:lba_full}
\small
\setlength{\tabcolsep}{3pt} % Tighten horizontal spacing for full 9-column width
\begin{tabular*}{\textwidth}{@{\extracolsep{\fill}}l l ccc ccc r}
\toprule
\textbf{Method} & \textbf{Model} & \multicolumn{3}{c}{\textbf{Seq. overlap 30\%}} & \multicolumn{3}{c}{\textbf{Seq. overlap 60\%}} & \textbf{Params} \\
\cmidrule(lr){3-5} \cmidrule(lr){6-8}
& & \textbf{RMSE} $\downarrow$ & \textbf{Pear.} $\uparrow$ & \textbf{Spear.} $\uparrow$ & \textbf{RMSE} $\downarrow$ & \textbf{Pear.} $\uparrow$ & \textbf{Spear.} $\uparrow$ & \\

\midrule
\textbf{No Pretraining} 
 & HoloProt-Full Surface \cite{Holoprot_somnath2021multi} & 1.464 & 0.509 & 0.500 & 1.365 & 0.749 & 0.742  & 1.4M\\
 & ProtNet-All-Atom \cite{ProtNet_wang2023}& 1.463 & \underline{0.551} & 0.551 & \underline{1.343} & \underline{0.765} & \underline{0.761}  & --- \\
 & ATOM3D-3DCNN \cite{LBA_dataset} & \underline{1.416} & 0.550 & \underline{0.553} & 1.621 & 0.608 & 0.615  & --- \\
 & ATOM3D-ENN \cite{LBA_dataset}& 1.568 & 0.389 & 0.408 & 1.620 & 0.623 & 0.633  & --- \\
 & ATOM3D-GNN \cite{LBA_dataset}& 1.601 & 0.545 & 0.533 & 1.408 & 0.743 &0.743  & ---\\
 & EPT-Scratch \cite{EPT_jiao2025} & \textbf{1.378} & \textbf{0.604}& \textbf{0.594}& \textbf{1.277}& \textbf{0.787}& \textbf{0.785} &30M\\
 & CT & 1.510& 0.501 & 0.486 & 1.386 & 0.736 & 0.720 &10M\\
 % & CGT & OOM & OOM & OOM & OOM & OOM & OOM &16M\\

\midrule
\textbf{Pretrained on PCQ} 
& Frad \cite{Frad_Ni2024}& 1.365& 0.599 & 0.577 & \textbf{1.213} & \textbf{0.804} & \textbf{0.801} &14M\\
& CT-SCD-PCQ & \textbf{1.332}& \textbf{0.617}& \textbf{0.600} & 1.226 & 0.800 & 0.794  &10M\\

\midrule
\textbf{Other Pretrained} 
 & EGNN-PLM \cite{EGNN-PLM_wu2023geometric} & 1.403 & 0.565 & 0.544 & 1.559 & 0.644 & 0.646  & 650M\\
 & Uni-Mol \cite{UniMol_Zhou2023} & 1.520 & 0.558 & 0.540 & 1.619 & 0.645 & 0.653  &47.6M\\
 & ProFSA \cite{profsa_2023} & 1.377 & 0.628 & 0.620 & 1.377 & 0.764 & 0.762  & 47.6M\\
 & EPT-Molecule \cite{EPT_jiao2025} & 1.336& 0.621& 0.602& 1.243& 0.802 & 0.800 &30M\\
 & EPT-Protein \cite{EPT_jiao2025} & 1.329& 0.628& 0.613& \underline{1.235} & \underline{0.804} & \underline{0.800} &30M\\
 & EPT-MultiDomain \cite{EPT_jiao2025}  & \underline{1.322} & \underline{0.644} & \underline{0.630} & 
 \textbf{1.227} & \textbf{0.811} & \textbf{0.803}  &30M\\
& ADiT-S \cite{lba-adit} & 1.337 & 0.626 & 0.618 & 1.413 & 0.740 & 0.740   & 12M \\
& ADiT-M \cite{lba-adit} & 1.353 & 0.622 & \underline{0.630} & 1.335  & 0.764 & 0.752  & 35M \\
& ADiT-L \cite{lba-adit} & \textbf{1.308} & \textbf{0.645} & \textbf{0.647} & 1.246  & 0.767 & 0.765  & 253M\\

\midrule
\textbf{SCD Pretrained} 
 & CT-SCD-AMP20 & 1.408 & 0.584 & 0.533 & 1.219 & 0.803 & 0.800 & 10M \\
 & CT-SCD-GEOM10 & 1.392 & 0.606 & 0.554 & 1.211 & 0.806 & 0.801 & 10M \\
 & CT-SCD-ALL & 1.372 & 0.594 & 0.578 & 1.218 & 0.802 & 0.798 & 10M \\
 & CT-SCD-SAIR-Lig & 1.412 & 0.566 & 0.547 & 1.283 & 0.779 & 0.770 & 10M \\
 & CT-SCD-SAIR & \underline{1.337} & \underline{0.617} & \underline{0.599} & \underline{1.196} & \underline{0.810} & 0.802 & 10M \\
 & CT-SCD-SAIR-Pocket & \textbf{1.304} & \textbf{0.640} & \textbf{0.624} & 1.200 & 0.809 & \underline{0.806} & 10M \\ 
 & CT-SCD-OMOL25 & 1.389 & 0.586 & 0.571 & \textbf{1.175} & \textbf{0.817} & \textbf{0.816} & 10M \\

\midrule
\textbf{FE Pretrained} 
& CT-FE-OMOL25 & 1.391 & 0.575 & 0.564 & 1.187 & 0.814 & 0.811 & 10M \\
\bottomrule
\end{tabular*}
\end{table*}

% \subsection{Matbench Properties: bandgap}
% % ######### Matbench: FULL Results ############### ---- Appendix
\begin{table*}[t]
\centering
\caption{\textbf{SCD on Periodic Materials}: Matbench band gap prediction performance for SCD and supervised force-energy (FE) pretrained models. Best results within each category are \textbf{bolded}; second best are \underline{underlined}. SCD models achieve competitive performance while utilizing significantly fewer parameters and unlabeled structures.}
\label{tab:mpgap_full}
\small
\begin{tabular*}{\textwidth}{@{\extracolsep{\fill}}l l c c r r}
\toprule
\textbf{Category} & \textbf{Model} & \multicolumn{2}{c}{\textbf{MP Gap (eV)}} & \textbf{\# Params} & \textbf{Pretraining} \\
\cmidrule(lr){3-4}
& & \textbf{Fold 0} & \textbf{Mean (0--4)} & & \textbf{Data Size} \\
\midrule
No Pretraining & MODNet \cite{MODNet_DeBreuck2021} & 0.215 & 0.220 & --- & n/a \\ 
 & coGN \cite{CoGN_ruff2023} & \textbf{0.153} & \textbf{0.156} & --- & n/a \\ 
 & JMP-S \cite{JMP_Shoghi2023} & 0.235 & --- & 25.2M & n/a \\ 
 & JMP-L \cite{JMP_Shoghi2023} & 0.228 & --- & 235M & n/a \\ 
 & CT (Baseline) & \underline{0.186} & --- & 10M & n/a \\
 % & CGT & nan & nan & 10M & n/a \\
\midrule
FE Pretrained & JMP-L \cite{JMP_Shoghi2023} & \textbf{0.089} & \textbf{0.091} & 235M & 120M \\ 
 & JMP-S \cite{JMP_Shoghi2023} & \underline{0.119} & \underline{0.121} & 27M & 120M \\
 & HackNIP (Orb-v2 + MODNet) \cite{HackNIP_kim2025} & --- & 0.150 & $>25.2$M & $>32$M\textsuperscript{*} \\
 & CT-FE-OMOL25 & 0.134 & --- & 10M & 4M \\

\midrule
Other SSL & Crystal-Twins \cite{xtal_twins} & --- &  0.264 & --- & 428k \\

\midrule
SCD Pretrained & CT-SCD-AMP20 & \textbf{0.122} & \textbf{0.123} & 10M & 675k \\
 & CT-SCD-ALL & \underline{0.132} & --- & 10M & 11.3M \\
 & CT-SCD-PCQ & 0.174 & --- & 10M & 3.4M \\
 & CT-SCD-GEOM10 & 0.177 & --- & 10M & 2.8M \\
 & CT-SCD-SAIR & 0.182 & --- & 10M & 4.4M \\
 & CT-SCD-OMOL25 & 0.136 & --- & 10M & 4M \\
\bottomrule
\end{tabular*}
\vspace{0.25em}
\flushleft\footnotesize \textsuperscript{*}Orb-v2 data size estimated from the sum of force-energy finetuning datasets: MPTraj (1.58M) and Alexandria (30.5M) \cite{HackNIP_kim2025}.
\end{table*}

\subsection{Force-Energy Pretraining and Fine-Tuning For Property Prediction}

In this work, we conduct all force-energy pretraining using the OMol25-4M dataset and a CT backbone. In Table \ref{tab:fe_pretrain}, we compare models pretrained via either direct ('forward') or gradient-based ('backward' aka 'conservative') force prediction. Unsurprisingly, the simple CT backbone does not yield state-of-the-art force-energy predictions; however, we find that SCD pretraining confers a modest benefit even when restricted to the same training geometries as the labeled task. Additionally, we find that a single SCD step (comprising a double forward pass) is roughly half the speed of a direct force training step (a single forward pass), but $1.4\times$ faster than a gradient-based (forward then 'backwards') force training update step. Table \ref{tab:fe_finetune} presents benchmarking results for these models on two representative QM9 tasks: HOMO and $U_0$ (potential energy at 0K). We observe that pretraining via direct force prediction benefits downstream property prediction more than gradient-based force pretraining, but it does not outperform SCD pretraining. In other tables, we report to our best model—pretrained exclusively on direct force prediction—as 'CT-FE-OMOL25'.

\begin{table}[h]
\centering
\caption{\textbf{OMol25 4M Force-Energy Pretraining}. Comparison of supervised force-energy training baselines (FE) against fine-tuning from SCD pretrained weights (SCD-FE). We compare direct force prediction (Forward) against gradient-based force calculation (Backward). CT-SCD-OMOL25 is pretrained by SCD using the same Omol25 4 geometries.}
\label{tab:fe_pretrain}
\small
\setlength{\tabcolsep}{5pt}
\begin{tabular}{l l l c cl}
\toprule
\textbf{Model Name} & \textbf{Initialization} & \textbf{Force Method} & \textbf{Val Force MAE} & \textbf{Val Energy MAE}  & \textbf{num training steps}\\
& & & (meV/\AA) & (meV)  &\\
\midrule
FE-Fwd & From Scratch & Forward & 44.8 & 648  &800k\\
FE-Fwd-long & From Scratch & Forward & 42.6 & 606  &1600k\\
FE-Bwd & From Scratch & Backward & 50.0 & 618  &800k\\

\midrule
SCD-FE-Fwd & CT-SCD-OMOL25 & Forward & 42.9 \small(+4.2\%) & \textbf{617} \small(+4.8\%)  &800k (+800k SCD)\\
SCD-FE-Bwd & CT-SCD-OMOL25 & Backward & \textbf{41.4} \small(+17.2\%) & 665 \small(-7.6\%) &800k (+800k SCD)\\
\bottomrule
\end{tabular}
\end{table}


\begin{table}[h]
\centering
\caption{\textbf{Fine-tuning Force-Energy Pretrained Models}. Force-energy pretrained models performed best when the energy head was reset and regularization noise was turned off during finetuning. Forward (non-conservative) force pretraining, conveyed the best average results on downstream property prediction and is reported as 'CT-FE-OMOL25' elsewhere.
}
\label{tab:fe_finetune}
\small
\setlength{\tabcolsep}{5pt}
\begin{tabular}{l l c cl}
\toprule
\textbf{Model} & \textbf{Fine-tuning Protocol} & \textbf{HOMO} & \textbf{$U_0$}  &avg \% impv. \\
& & (meV) & (meV)  &over baseline\\
\midrule
 baseline& & 20.3& 6.15&---\\
\midrule
FE-Fwd & Standard & 15.9 & 12.9  &-44.0\%\\
FE-Fwd & Reset Head, $\sigma_{\text{reg}}=0.005$ & 15.0 & 10.2  &-19.9\%\\
FE-Fwd & Reset Head, $\sigma_{\text{reg}}=0$ & 13.6 & 3.40  &38.9\%\\
FE-Fwd-long & Reset Head, $\sigma_{\text{reg}}=0$ & 13.7 & 3.47  &38.0\%\\
FE-Bwd & Reset Head, $\sigma_{\text{reg}}=0$ & 14.5 & 3.36  &37.0\%\\
\midrule
CT-SCD-OMol25 & Standard & 12.8 & 3.54  & 39.7\%\\
\bottomrule
\end{tabular}
\end{table}





% \subsection{QM9 results from other SSL studies using different architectures and pretraining datasets}

\begin{table}[h]
\centering
\caption{SCD vs. other SSL Methods using different backbones and pretraining data. Models CT and CGT below are both pretrained on the PCQ dataset. Best results are \textbf{bolded}; second best are \underline{underlined}.}
\label{tab:scd_vs_SSL_qm9}
\small
\setlength{\tabcolsep}{3pt}
\begin{tabular}{l c c c c}
\toprule
\textbf{QM9 Task} & \textbf{CT} & \textbf{CGT} & \textbf{EPT-10} & \textbf{Uni-Corn} \\
 & (Ours) & (Ours) & \cite{EPT_jiao2025} & \cite{UniCorn_Feng2024} \\
\midrule
Params & 10M & 17M & 30M & --- \\
Pretrain Data & 3.4M & 3.4M & 5.9M & 15M \\
\midrule
HOMO (meV) & \underline{12.7} & \textbf{9.65} & 15.2 & 13.0 \\
LUMO (meV) & \underline{11.5} & \textbf{9.05} & 13.6 & 11.9 \\
Gap (meV) & \underline{24.5} & \textbf{19.7} & 29.0 & 24.9 \\
ZPVE (meV) & \underline{1.18} & 1.22 & \textbf{1.11} & 1.4 \\
$U_0$ (meV) & \textbf{3.58} & \underline{3.96} & 5.44 & 3.99 \\
$U$ (meV) & \textbf{3.50} & 4.11 & 5.54 & \underline{3.95} \\
$H$ (meV) & \textbf{3.52} & 4.00 & 5.42 & \underline{3.94} \\
$G$ (meV) & \underline{5.29} & \underline{5.29} & 6.37 & \textbf{5.09} \\
$\alpha$ ($a_0^3$) & \underline{0.0377} & 0.0383 & 0.045 & \textbf{0.036} \\
$C_v$ (cal/mol K) & 0.021 & \textbf{0.019} & 0.020 & \textbf{0.019} \\
\bottomrule
\end{tabular}
\vspace{0.25em}
\flushleft\footnotesize The pretraining datasets for EPT-10 and Uni-Corn include the 3.4M PCQ structures used by SCD.
\end{table}


\subsection{MD17: Small Molecule Forces}

Consistent with prior work, we benchmark a smaller SCD model (6 layers, dim 128) on MD17 \cite{md17_chmiela2017machine}. The models presented in Table \ref{tab:md17_scd} were pre-trained on the PCQ dataset and subsequently fine-tuned on MD17. To match the fine-tuning methods described by Frad \cite{Frad_Ni2024}, we also train a model using a two pass 'force-denoise' scheme, in which the first pass is used to predict forces and energies and the second forward pass is used for a auxiliary denoising task. We observe that SCD pre-training with the CT backbone yielded results similar to those of Frad and SLiDe, while achieving 4x faster inference speeds. However, none of these methods achieve particularly strong performance on this benchmark. We hypothesize that this stems from the pre-training dataset containing exclusively ground-state geometries; consequently, the non-equilibrium structures in MD17 are likely out-of-distribution. Lacking non-equilibrium structures during pre-training, the models fail to capture features necessary for non-equilibrium prediction tasks. This aligns with our findings on other benchmarks, where similarity between pre-training and fine-tuning data distributions is a strong predictor of performance. We do not report results for the GET/CGT variant due to persistent numerical instability (segmentation faults) when computing gradient-based forces.

% Regardless, GET is significantly slower than ET and CT, making it poorly suited to downstream force/energy prediction applications.]
\begin{table}[t]
\centering
\caption{SCD achieves comparable performance on MD17 force--energy prediction to prior methods while using a smaller and faster architecture. Best results are in \textbf{bold}, second best are \underline{underlined}.}
\label{tab:md17_scd}
\renewcommand{\arraystretch}{1.15}
\setlength{\tabcolsep}{6pt}
\resizebox{\linewidth}{!}{%
\begin{tabular}{l|cccc!{\vrule width 1.2pt}cc}
\toprule
\multicolumn{1}{l|}{\textbf{Method}}              
& \textbf{Baseline} 
& \textbf{Coord} 
& \textbf{SCD} 
& \textbf{SCD (force-denoise)} 
& \textbf{Frad (force-denoise)} 
& \textbf{SliDe} \\

\multicolumn{1}{l|}{\textbf{Architecture}}        
& ET-small 
& ET-small 
& CT-small 
& CT-small 
& GET-small 
& GET-small \\

\multicolumn{1}{l|}{\textbf{Pretraining dataset}} 
& none 
& PCQ 
& PCQ 
& PCQ 
& PCQ 
& PCQ \\
\midrule
\multicolumn{7}{l}{\textbf{Forces MAE (kcal/mol/\AA)}}\\
\midrule
Aspirin         & 0.253 & 0.211 & 0.210 & \underline{0.200} & 0.209 & \textbf{0.174} \\
Benzene        & 0.196 & \textbf{0.169} & 0.171 & 0.175 & 0.199 & \textbf{0.169} \\
Ethanol        & 0.109 & 0.096 & 0.120 & 0.116 & \underline{0.091} & \textbf{0.088} \\
Malonaldehyde  & 0.169 & 0.139 & 0.186 & 0.173 & \textbf{0.142} & 0.154 \\
Naphthalene    & 0.061 & 0.053 & \underline{0.041} & \textbf{0.040} & 0.053 & 0.048 \\
Salicylic acid & 0.129 & 0.109 & \underline{0.101} & \textbf{0.099} & 0.108 & \underline{0.101} \\
Toluene        & 0.067 & 0.058 & \underline{0.053} & \textbf{0.051} & 0.054 & 0.054 \\
Uracil         & 0.095 & 0.074 & 0.077 & \textbf{0.072} & \underline{0.076} & 0.083 \\
\midrule[1.2pt]
\multicolumn{7}{l}{\textbf{Energy MAE (kcal/mol)}}\\
\midrule
Aspirin         & 0.123 & -- & 0.0742 & \textbf{0.0705} & -- & -- \\
Benzene         & 0.058 & -- & 0.0274 & \textbf{0.0211} & -- & -- \\
Ethanol         & 0.052 & -- & 0.0286 & \textbf{0.0280} & -- & -- \\
Malonaldehyde   & 0.077 & -- & 0.0432 & \textbf{0.0417} & -- & -- \\
Naphthalene     & 0.085 & -- & 0.0163 & \textbf{0.0156} & -- & -- \\
Salicylic acid  & 0.093 & -- & \textbf{0.0286} & 0.0290 & -- & -- \\
Toluene         & 0.074 & -- & \textbf{0.0181} & 0.0184 & -- & -- \\
Uracil          & 0.090 & -- & \textbf{0.0190} & 0.0193 & -- & -- \\
\bottomrule
\end{tabular}
}
\vspace{0.5em}
\caption*{\footnotesize CGT results on MD17 are omitted due to an unresolved segmentation fault during backpropagation.}
\end{table}


\section{Datasets}
\label{appendix_datasets}

\subsection{Pretraining Datasets}
% We've made all pretraining datasets available for download on huggingface here: (add link) 

\textbf{PCQ}: The PCQM4Mv2 dataset is a curated subsample of the PubChemQC project containing ground state geometries for 3.378M organic small molecule obtained using B3LYP/6-31G* DFT \cite{PubChemQC_Nakata2017}. This dataset is commonly used for demonstrating self-supervised learning methods, either on its own or as part of a composite dataset. Throughout this work we use 'PCQ' to indicate models trained only on this dataset. 

% We've uploaded a copy of this preprocessed dataset to huggingface here: (add link)

\textbf{GEOM10}: The original GEOM dataset contains 36M conformers from $\sim$450k unique molecules generated using GFN2-xTB, a semi-empirical extended tight-binding method \cite{GEOM_axelrod2022}. GEOM is split into three categories: 1) GEOM-drugs, containing conformers of drug like small molecules, 2) GEOM-QM9, conformers of molecules from the qm9 dataset, and 3) GEOM-MoleculNet, containing conformers from a subsample of molecules in the MoleculeNet benchmark. To avoid overlap with benchmarking datasets we only utilize structures from geom-drug, containing 304k unique molecules, and curate two subsamples: GEOM10, containing up to 10 of the lowest energy conformers from each molecule (2.7M geometries total), and GEOM1 containing only the lowest energy (ground-state) conformer from each molecule (304k geometries). 
% We've made our versions of GEOM available on huggingface here: (add link)

\textbf{SAIR}: The Structurally Augmented IC50 Repository (SAIR) dataset contains a large number of bound protein-ligand structures generated using the Boltz-1x model \cite{boltz1_2024}. Five conformers are provided for each structure. Although, the authors of the SAIR reports to have 1,048,857 unique binding pairs totaling in 5.2M conformers, the copy we obtained from huggingface (https://huggingface.co/datasets/SandboxAQ/SAIR) contains only 888,104 unique binding pairs and 4.4M conformers. 

\textbf{Alex-MP-20}: The Alex-MP-20 dataset comprises 607,673 periodic materials structures with up to 20 atoms \cite{mattergen_zeni2023} derived from both the Materials Project \cite{MaterialsProject}, and the Alexandria materials database \cite{Alexandria_1, Alexandria_2}. Throughout this work, we refer to Alex-MP-20 as 'AMP20'. When SCD pretraining on AMP20, we apply a unit cell repeat augmentation to simulate larger crystals. Each sample is randomly tiled in one direction with a probability ($p_{\text{cell repeat}}$) of 50\%, for up to two repeats. This dataset is also available on huggingface (https://huggingface.co/datasets/OMatG/Alex-MP-20).

\textbf{AllAtoms (ALL):} We also pretrain on a combine dataset that includes the entirety of PCQ, GEOM10, SAIR, and AlexMP20, totaling in 5.2M unique structures with 11.3M conformations. 


\textbf{OMol25 4M split}: The Open Molecular Dataset 2025 (OMol25) is a large-scale repository of over 100 million quantum chemistry calculations covering 83 million unique molecular systems \cite{OMOL25_levine2025open}. While the full dataset spans the first 83 elements (H--Bi) and systems up to 350 atoms, we utilize the designated `4M split', a randomly sampled subset of approximately 4 million geometries. All labels (forces and energies) were computed at the $\omega$B97M-V/def2-TZVPD level of theory. The 4M subset preserves the broad chemical diversity of the full dataset, containing equilibrium and non-equilibrium structures from four primary domains: small organic molecules, biomolecules (proteins, DNA/RNA from PDB/BioLiP2), metal complexes (transition metals, ligands), and electrolytes. This dataset is available on huggingface (https://huggingface.co/facebook/OMol25).


\subsection{Benchmarking Datasets and Finetuning Hyperparameters}

\textbf{QM9}: The QM9 dataset is a commonly used benchmark for small molecule property prediction on 3D geometries\cite{QM9_Ramakrishnan2014}. QM9 comprises a set of small organic molecules containing up to nine heavy atoms (C,N,F,O), or up to 29 atoms total when including hydrogen. All 3D Geometries from QM9 are computed using the B3LYP/6-31G(2df, p) DFT functional and are labeled with a set of 15 quantum mechanical properties including homo-lumo gap, atomization energies, specific heat, polarizability and more. The original dataset contains 134k molecules, however it is now standard practice to use a cleaned version containing 130,831. We obtain our copy from the PyTorch Geometric library. QM9 is typically split randomly: 110,000 for training, 10,000 for validation, and the remaining 10,831 for the test set. We use this convention in our work, however we do not use the validation set for early stopping. We follow the practices described in TorchNet-MD\cite{TorchMD-Net_thölke2022} and subtract provided references energies from atomization energy targets. All other targets are normalized to a standard gaussian for prediction.


\textbf{LBA}: The Ligand Binding Affinity (LBA) dataset comprises 4,463  co-crystallized 3D structures of  protein-ligand complexes derived from PDBbind. Binding affinities targets are experimentally derived and provided as negative log of either inhibition or dissociation constant ($-log(K)$). Two splits are provided based on sequence overlap between proteins in the training set and test set; in ‘id60’ no protein in the training set has more than a 60\% sequence overlap with any protein in the testset. Similarly, in ‘id30’ train-test sequence overlap is limited to 30\%. Following ATOM3D preprocessing\cite{LBA_dataset}, predictions are made using extracted pocket-ligand complexes.


\textbf{Matbench}: The matbench properties dataset \cite{matbench_2020} provides 13 different tasks (regression and classification) for periodic materials structures to benchmark property prediction models. Across these tasks, training set size ranges from 312-132k samples, each of which are split into 5 folds with predefined test splits. In this work we explore only one regression task, 'mp\_gap' or bandgap energy prediction, as a proof of concept. We choose bandgap energy because of its economic relevance. The ‘mp\_gap’ task provides 106k periodic crystal structures in total, resulting in training set sizes of 84k and test set sizes of 21k. We report results for split 0 with all models tested. For our best model we also report the average of all splits.

\textbf{MD17}: The MD17 \cite{md17_chmiela2017machine} dataset comprises a collection of ab initio molecular dynamics (AIMD) trajectories for eight small organic molecules with DFT forces and energies, commonly used to evaluate MLIPs. Following past convention we use a random split of 950 samples for training, 50 for validation (however we do not employ early stopping), and the remainder as the test set. The total number of geometries provided for each trajectory varies from 150k-1M. All force predictions provided are derived from energy gradients (a backwards pass).



\section{Hyperparameters}






\subsection{Constants across all training runs}
% \FloatBarrier

All training runs discussed in this work use the AdamW optimizer with $\beta_2$=0.999, and a learning rate scheduler composed of a linear warm up, followed by cosine decay. No validation sets are used for early stopping, as is often done in other works. All graph creation uses a cutoff distance of 5.0 \AA.

% \subsection{Pretraining}

\begin{table}[t]
\centering
\caption{Base SCD pretraining hyperparameters: applied in all SCD training runs }
\label{tab:pretrain_hparams}
\begin{tabular}{lc}
\toprule
\textbf{Hyperparameter} & \textbf{Value} \\
\midrule
Warmup steps        & 10{,}000 \\
Learning rate       & 0.005 \\
$\beta_1$           & 0.9 \\
Weight decay        & 0.05 \\
DropPath            & 0.1 \\
Corrupt noise $\sigma$ & 0.04 \\
Regularization noise $\sigma$ & 0.005 \\
\bottomrule
\end{tabular}
\end{table}

\begin{table}[t]
\centering
\caption{SCD Pretraining: dataset and model specific hyperparameters.}
\label{tab:scd_pretrain_hparams}
% % \begin{adjustbox}{max width=\linewidth}
\begin{tabular}{lccccccc}
\toprule
% \textbf{Name} &
% \textbf{CT-SCD-PCQ} &
% \textbf{CGT-SCD-PCQ} &
% \textbf{CT-SCD-GEOM10} &
% \textbf{CT-SCD-AMP20} &
% \textbf{CT-SCD-SAIR} &
% \textbf{CT-SCD-ALL} &
% \textbf{CT-SCD-Omol25} \\
\midrule
Backbone            & CT  & CGT & CT & CT & CT & CT & CT \\
Dataset             & PCQ & PCQ & GEOM10 & AMP20 & SAIR & ALL & Omol25 \\
Total steps         & 800k & 800k & 800k & 800k & 800k & 1.2M & 800k \\
Batch size / GPU    & 128 & 32 & 48 & 142 & 12 & 54 & 48 \\
\# GPUs (A100)      & 4 & 16 & 16 & 4 & 8 & 16 & 16 \\
Effective batch size& 512 & 512 & 768 & 568 & 96 & 864 & 768 \\
\midrule
\multicolumn{8}{l}{\textbf{Materials augmentations}} \\
Cell repeat iters   & -- & -- & -- & 2 & -- & 2 & -- \\
$p_{\text{cell repeat}}$ & -- & -- & -- & 0.5 & -- & 0.5 & -- \\
\bottomrule
\end{tabular}
% \end{adjustbox}
\end{table}

% \FloatBarrier

% \subsection{Finetuning}
% \subsubsection{QM9}
% For QM9 we use a common data split of 110,000 training samples

% train 110000, val 10000, test 10831


\begin{table}[t]
\centering
\caption{QM9 finetuning hyperparameters for different target groups.}
\label{tab:qm9_finetune_hparams}
% \begin{adjustbox}{max width=\linewidth}
\begin{tabular}{lccc}
\toprule
\textbf{Hyperparameter} &
\textbf{HOMO, LUMO, Gap} &
\textbf{U0, U, H, G} &
\textbf{Alpha, ZPVE, CV} \\
\midrule
Batch size                 & 128   & 32    & 128   \\
Learning rate              & 5e{-4} & 3e{-4} & 3e{-4} \\
Total steps                & 300k  & 500k  & 400k  \\
Warmup steps               & 10k   & 10k   & 10k   \\
$\beta_1$                  & 0.995 & 0.995 & 0.95  \\
Weight decay               & 0.01  & 0.01  & 0.01  \\
DropPath (init / final)    & 0.15 / 0.1 & 0.15 / 0.0 & 0.15 / 0.01 \\
Reg. noise $\sigma$        & 0.005 & 0.0 & 0.005 \\
Denoise loss weight        & 0.1   & 0.0   & 0.2   \\
Subtract ref. energy       & No    & Yes   & No    \\
Loss EMA                   & Off   & 0.05  & Off   \\
Head aggregation           & Sum   & Sum   & Sum   \\
\bottomrule
\end{tabular}
% \end{adjustbox}
\end{table}


% \subsubsection{Matbench: materials project band gap}

\begin{table}[t]
\centering
\caption{Matbench finetuning hyperparameters for the \texttt{mp\_gap} task.}
\label{tab:matbench_mpgap_hparams}
\begin{tabular}{lc}
\toprule
\textbf{Hyperparameter} & \textbf{Value} \\
\midrule
Batch size                 & 8 \\
Learning rate (from scratch) & 2e{-4} \\
Learning rate (finetuning)   & 8e{-5} \\
Total steps                & 300k \\
Warmup steps               & 10k \\
$\beta_1$                  & 0.995 \\
Weight decay               & 0.01 \\
DropPath (init / final)    & 0.15 / 0.00 \\
Reg. noise $\sigma$        & 5e{-4} \\
Denoise loss weight        & 0.01 \\
Head aggregation           & Sum \\
\bottomrule
\end{tabular}
\end{table}

% \subsubsection{Ligand Binding Affinity}
% Note details on: conditional denoising lig vs pocket. Atom masking in output

\begin{table}[t]
\centering
\caption{Finetuning hyperparameters for the Ligand Binding Affinity (LBA) benchmark}
\label{tab:lba_hparams}
\begin{tabular}{lc}
\toprule
\textbf{Hyperparameter} & \textbf{Value} \\
\midrule
Batch size                  & 8 \\
Learning rate (from scratch) & 2e{-4} \\
Learning rate (finetuning)   & 1e{-4} \\
Total steps                 & 30k \\
Warmup steps                & 500 \\
$\beta_1$                   & 0.995 \\
Weight decay                & 0.05 \\
DropPath                    & 0.15 \\
Reg. noise $\sigma$         & 0.005 \\
Denoise loss weight         & 0.1 \\
Head aggregation            & Mean \\
\bottomrule
\end{tabular}
\end{table}


% \subsubsection{MD17}
\begin{table}[t]
\centering
\caption{Training hyperparameters for the MD17 benchmark.}
\label{tab:md17_hparams}
\begin{tabular}{lc}
\toprule
\textbf{Hyperparameter} & \textbf{Value} \\
\midrule
Backbone                  & CT-small \\
Training set size         & 1{,}000 \\
Batch size                & 8 \\
Energy loss EMA           & 0.05 \\
Energy loss weight        & 0.8 \\
Force loss weight         & 0.2 \\
Learning rate             & 5e{-4} \\
Warmup steps              & 5k \\
Total steps               & 150k \\
$\beta_1$                 & 0.995 \\
Weight decay              & 0.01 \\
DropPath (init / final)   & 0.15 / 0.00 \\
Reg. noise $\sigma$       & 0.005 \\
Denoise loss weight       & 0.1 \\
Head aggregation          & Sum \\
\bottomrule
\end{tabular}
\end{table}




%%%%%%%%%%%%%%%%%%%%%%%%%%%%%%%%%%%%%%%%%%%%%%%%%%%%%%%%%%%%%%%%%%%%%%%%%%%%%%%
%%%%%%%%%%%%%%%%%%%%%%%%%%%%%%%%%%%%%%%%%%%%%%%%%%%%%%%%%%%%%%%%%%%%%%%%%%%%%%%

\end{document}

% This document was modified from the file originally made available by
% Pat Langley and Andrea Danyluk for ICML-2K. This version was created
% by Iain Murray in 2018, and modified by Alexandre Bouchard in
% 2019 and 2021 and by Csaba Szepesvari, Gang Niu and Sivan Sabato in 2022.
% Modified again in 2023 and 2024 by Sivan Sabato and Jonathan Scarlett.
% Previous contributors include Dan Roy, Lise Getoor and Tobias
% Scheffer, which was slightly modified from the 2010 version by
% Thorsten Joachims & Johannes Fuernkranz, slightly modified from the
% 2009 version by Kiri Wagstaff and Sam Roweis's 2008 version, which is
% slightly modified from Prasad Tadepalli's 2007 version which is a
% lightly changed version of the previous year's version by Andrew
% Moore, which was in turn edited from those of Kristian Kersting and
% Codrina Lauth. Alex Smola contributed to the algorithmic style files.
